% Vorlage fuer Abschlussarbeiten am Fachbereich Mathematik der TU Darmstadt.
% Geeignet fuer Bachelorarbeiten, Masterarbeiten und Diplomarbeiten.

%\documentclass[11pt,a4paper,twoside,openright]{report}
\documentclass[12pt,a4paper,twoside]{report}
% Benutzt man die Option draft, so kann man die Umbrueche ueberpruefen.

% Hier werden alle benoetigten Pakete und Einstellungen geladen. Auch
% hier sind Sie frei diese direkt in der Praeambel zu laden.


\usepackage[ngerman]{babel}

% Kodierung fuer unixoide Systeme und Windowssysteme.

\usepackage[utf8]{inputenc}

% Schriftart Times New Roman.

\usepackage{times}

% Erweitert den Zeichenvorrat, so dass z.B. auch Umlaute im PDF-Dokument
% gefunden werden.

\usepackage[T1]{fontenc}

% Zur Einbindung von Bildern.

\usepackage{graphicx}

% Erweiterte enumerate-Umgebung.

\usepackage{enumerate}

% Verschiedene Pakete, die nuetzlich sind um Mathematik in LaTeX zu setzen.

\usepackage{amsmath, amssymb, amsthm, dsfont}

% Indentation

\usepackage{scrextend}

% Spaces for newcommand

\usepackage{xspace}

% tikz
\usepackage{tikz}
%\usetikzlibrary{arrows.meta}% arrow tip library
%\usetikzlibrary{bending}% better arrow head for bended lines

% Hyphenation of wordas that already contain a hyphen
% http://tex.stackexchange.com/questions/2706/adequate-hyphenation-of-words-already-containing-a-hyphen
% Example:
% \textsc{Alexandroff}\hyp{}\textsc{Urysohn}\hyp{}Kompaktheit

\usepackage{hyphenat}

% Hyperlinks

%\usepackage{hyperref}
\usepackage[pdfauthor={Fabian Gabel},%
pdftitle={Bachelorthesis - Helmholtz-Zerlegung},%
%pagebackref=true,%
%pdftex
]{hyperref}


\usepackage[
%nonumberlist, %keine Seitenzahlen anzeigen
toc,          %Einträge im Inhaltsverzeichnis
section]      %im Inhaltsverzeichnis auf section-Ebene erscheinen
{glossaries}

%%%%%%%%%%%%%%%%
% Seitenlayout %
%%%%%%%%%%%%%%%%

% DIV# gibt den Divisor für die Layoutberechnung an.
% Vergrößern des Divisors vergrößert den Textbereich.
% BCOR#cm gibt die Breite des Bundstegs an.
\usepackage[DIV14,BCOR2cm]{typearea}

% Abstand obere Blattkante zur Kopfzeile ist 2.54cm - 15mm
\setlength{\topmargin}{-15mm}

% Keine Einrueckung nach einem Absatz.

\parindent 0pt

% Abstand zwischen zwei Abs\"atzen.

\parskip 12pt

% Zeilenabstand.

\linespread{1.25}

% Inhaltsverzeichnis erstellen.

\usepackage{makeidx}
\makeindex


% Hier werden Makros und Umgebungen eingebunden. Diese werden separat
% in der Datei befehle.tex definiert. Sie sind frei diese Befehle auch
% direkt in der Praeambel zu definieren.
% Makros

\newcommand{\N}{\mathbb{N}} % natuerliche Zahlen
\newcommand{\Z}{\mathbb{Z}} % ganze Zahlen
\newcommand{\Q}{\mathbb{Q}} % rationale Zahlen
\newcommand{\R}{\mathbb{R}} % reelle Zahlen
\newcommand{\K}{\mathbb{K}} % Körper
\newcommand{\co}{\mathrm{co}} % compact open
\newcommand{\W}{\mathcal{W}} % Weyl-Gruppe

\DeclareMathOperator{\Ret}{Re} % Realteil
%\newcommand{\PIT}{\textbf{PIT }} % Prime Ideal Theorem
\newcommand{\AC}{\textbf{AC}\xspace} % Prime Ideal Theorem
\newcommand{\UFT}{\textbf{UFT}\xspace} % Prime Ideal Theorem
\newcommand{\PIT}{\textbf{PIT}\xspace}
\newcommand{\ZF}{\textbf{ZF}\xspace}
\newcommand{\ZFC}{\textbf{ZFC}\xspace}
\newcommand{\CC}{\textbf{CC}\xspace}
\newcommand{\CCR}{\textbf{CC}($\R$)\xspace}
\newcommand{\PCC}{\textbf{PCC}\xspace}
\newcommand{\PCCR}{\textbf{PCC}($\R$)\xspace}

\newcommand{\grad}[1]{\gv{\nabla} #1} % for gradient
\let\divsymb=\div % rename builtin command \div to \divsymb
\renewcommand{\div}[1]{\mathrm{div} #1} % for divergence
\renewcommand{\d}[1]{\ensuremath\, {\operatorname{d}\!{#1}}}

\newcommand{\banach}{\textsc{Banach}}
\newcommand{\cauchy}{\textsc{Cauchy}}
\newcommand{\euklid}{\textsc{Euklid}}
\newcommand{\fubini}{\textsc{Fubini}}
\newcommand{\hahn}{\textsc{Hahn}}
\newcommand{\helmholtz}{\textsc{Helmholtz}}
\newcommand{\hoelder}{\textsc{Hölder}}
\newcommand{\killing}{\textsc{Killing}}
\newcommand{\lebesgue}{\textsc{Lebesgue}}
\newcommand{\leibniz}{\textsc{Leibniz}}
\newcommand{\lie}{\textsc{Lie}}
\newcommand{\lipschitz}{\textsc{Lipschitz}}
\newcommand{\poincare}{\textsc{Poincaré}}
\newcommand{\sobolev}{\textsc{Sobolev}}
\newcommand{\weyl}{\textsc{Weyl}}


% Umgebungen für Definitionen, Sätze, usw.

%\newtheorem{defn}{Definition}[section]
%\newtheorem{ex}{Beispiel}[chapter]

\theoremstyle{plain}
\newtheorem{thm}{Satz}[chapter]
\newtheorem{lem}[thm]{Lemma}
\newtheorem{cor}[thm]{Korollar}
\newtheorem{prop}[thm]{Proposition}

\theoremstyle{definition}
\newtheorem{defn}[thm]{Definition}

\theoremstyle{remark}
\newtheorem*{bem}{Bemerkung}

\def\Satzrefname{Satz}

\DeclareMathOperator{\spn}{span}
\DeclareMathOperator{\sign}{sign}
\DeclareMathOperator{\GL}{GL}
\DeclareMathOperator{\Sym}{Sym}
\DeclareMathOperator{\hgt}{ht}
\DeclareMathOperator{\dist}{dist}
\DeclareMathOperator{\loc}{loc}

%\renewcommand{\refname}{Literaturverzeichnis}
\addto\captionsngerman{%
\renewcommand*{\bibname}{Literaturverzeichnis}
}

%\makeglossaries

\pagestyle{fancy}
\fancypagestyle{MyStyle}{%
\fancyhf{} %Clean headers
\fancyhead[LE,RO]{\thepage}
\fancyhead[RE]{\normalfont \nouppercase{\leftmark}}
\fancyhead[LO]{\normalfont \nouppercase{\rightmark}}
\fancyfoot[C]{\tiny Version vom \today}
\renewcommand{\chaptermark}[1]{\markboth{\thechapter\ {\normalfont{##1}}}{}}
\renewcommand{\sectionmark}[1]{\markright{\thesection\ {\normalfont{##1}}}{}}
}

\fancypagestyle{plain}{%
\fancyhf{} % clear all header and footer fields
\fancyfoot[LE,RO]{\thepage} % except the center
\renewcommand{\headrulewidth}{0pt}
\renewcommand{\footrulewidth}{0pt}}


\begin{document}
% Auf der Titelseite und im Inhaltsverzeichnis sollen keine
% Seitenzahlen dargestellt werden.
\pagestyle{empty}

% Hier wird die Titelseite eingebunden.
\begin{titlepage}
  \begin{center}
    \vspace{1cm}
    \includegraphics[width=0.5\linewidth]{TU_Darmstadt_Logo.pdf}
    \vspace{1cm}
    
    \large{Fachbereich Mathematik}
    \vspace{2.5cm}
    
    \large{Bachelorarbeit}
    \vspace{2cm}

    \Huge{Die Helmholtz-Zerlegung in $L^2$}
    
    \vspace*{3cm}    
    
		\large
                \href{mailto:gabel@mathematik.tu-darmstadt.de}{Fabian Gabel}
    \vspace*{1.0cm}

    15.10.2016 \\
    \vspace*{2cm}

    Betreuer: PD Dr. Robert Haller-Dintelmann

    \vspace*{.5cm}

    %Zweiter Gutachter: Name des zweiten Gutachters\\[2ex]
    \vspace*{\fill}
    \tiny{Version vom \today}
  \end{center}
\end{titlepage}
\vspace*{\fill}


% Inhaltsverzeichnis erstellen.
\tableofcontents

% Ab sofort werden Seitenzahlen in der Kopfzeile angezeigt.
%\pagestyle{headings}
\pagestyle{MyStyle}

\chapter*{Einleitung}\index{Einleitung}
\addcontentsline{toc}{chapter}{Einleitung}\index{Einleitung}

\begin{addmargin}[2em]{2em}% 1em left, 2em right
  \textit{... as Sir Cyril Hinshelwood has observed ... fluid dynamicists
were divided into hydraulic engineers who observed things that
could not be explained and mathematicians who explained things
that could not be observed.} 
  \flushright(James Lighthill)
\end{addmargin}
\vspace{1.5cm}

Die akkurate Modellierung des Verhaltens \newton scher Fluide ist zentral für unzählige Anwendungen der Aerodynamik, Verbrennungsforschung oder chemischen Industrie.
Die Grundlage für die Modellierung bildet ein System partieller Differentialgleichungen, welches \navier\ und \stokes\ unabhängig voneinander in der ersten Hälfte des 19. Jahrhunderts einführten. 
Man bezeichnet dieses Gleichungssystem heute als \navier\hyp\stokes\hyp{}Gleichungen.
Sie umfassen im inkompressiblen dreidimensionalen Fall eine Gleichung zur Beschreibung der Massenerhaltung, die sogenannte Kontinuitätsgleichung, und für jede Raumrichtung eine Impulsgleichung:

\begin{align*}
  \div u &= 0, \\
  \frac{\partial u}{\partial t} - \nu \Delta u + u \cdot \nabla u + \nabla p &= f.
\end{align*}

Die Existenz und Eindeutigkeit klassischer glatter Lösungen dieses Gleichungssystems gehört nach dem \textsc{Clay}\hyp{}Institute zu einem der wichtigsten ungelösten mathematischen Probleme unseres Jahrtausends und ist daher in der Liste der Millennium Probleme zu finden.

In der Lösungstheorie der \navier\hyp\stokes\hyp{}Gleichungen aber auch partieller Differentialgleichungen im Allgemeinen hat es sich als nützlich erwiesen, auf der Suche nach klassischen Lösungen einen Umweg einzuschlagen und zunächst die Existenz sogenannter schwacher Lösungen nachzuweisen.
In einem zweiten Schritt wird dann überprüft, ob die gefundenen schwachen Lösungen über zusätzliche Regularitätseigenschaften verfügen, welche die Existenz einer Lösung im klassischen Sinne garantieren.

Als nützliche Werkzeuge für die Lösungstheorie haben sich dabei Zerlegungen geeigneter Funktionenräume erwiesen. Die Idee dazu fußt auf einer von \helmholtz\ erstmals 1870 formulierten Zerlegung eines Vektorfeldes in ein skalares Potential und ein Vektorpotential.
Im Rahmen der Suche nach schwachen Lösungen im \hilbert\hyp{}Raum $L^2(\Omega)^n$, wobei $\Omega$ ein Teilgebiet des $\R^n$, $n \geq 2$, bezeichne, findet auch das in dieser Arbeit vorgestellte Hilfsmittel Verwendung: die \helmholtz\hyp{}Zerlegung des $L^2$.
Dabei handelt es sich um eine orthogonale Zerlegung des Lösungsraumes, welche es unter anderem ermöglicht den neben dem gesuchten Geschwindigkeitsfeld $u$ unbekannten Druck $p$ vorerst aus dem Gleichungssystem zu eliminieren und so die Anzahl der Unbekannten zu reduzieren.
Die Anwendung der \helmholtz\hyp{}Zerlegung auf ein Element des Lösungsraumes zerlegt dieses in einen divergenzfreien und einen rotationsfreien Anteil.

Ziel dieser Arbeit ist es, aufbauend auf den im Grundstudium vermittelten Kenntnissen der Funktionentheorie, Integrationstheorie und Funktionalanalysis, die Existenz und Eindeutigkeit der \helmholtz\hyp{}Zerlegung auf dem Raum $L^2(\Omega)^n$ zu beweisen.

Das erste Kapitel stellt dazu die nötigen funktionalanalytischen Grundlagen bereit und führt in die in der Arbeit verwendete Notation ein.
Das Hauptaugenmerk dieses Kapitels liegt auf der Definition der schwachen Differenzierbarkeit und des Distributionsbegriffs.

Im zweiten Kapitel werden die zum Beweis der \helmholtz\hyp{}Zerlegung nötigen Hilfsaussagen bereitgestellt.
Das zentrale Resultat dieses Teils der Arbeit ist ein Kriterium, welches unter gewissen Bedingungen die schwache Lösbarkeit der Gradientengleichung $\nabla p = f$ sicherstellt.

Kapitel drei behandelt schließlich die \helmholtz\hyp{}Zerlegung $$L^2(\Omega)^n = L^2_\sigma(\Omega) \oplus G_2(\Omega)$$ und eine Charakterisierung dieser Zerlegung für den Fall, dass $\Omega = \R^n$ gilt.

Das letzte Kapitel fasst nochmals die zentralen Resultate zusammen und gibt einen Ausblick auf Problemstellungen, in denen die \helmholtz\hyp{}Zerlegung zum Einsatz kommt.

\chapter{Funktionalanalytische Grundlagen}
\label{cp:grundlagen}

Dieses Kapitel beschäftigt sich mit Definitionen und Eigenschaften der für die kommenden Kapitel zentralen Funktionenräume.
Die Notation orientiert sich an \cite{sohr2001navier}.
In Anlehnung an die für \hilbert\hyp{}Räume verbreitete Schreibweise $\langle \cdot, \cdot \rangle$ für das Skalarprodukt wollen wir nicht nur für den Raum $L^2$, sondern auch für alle anderen $L^q$\hyp{}Räume diese Symbolik verwenden.
Wir definieren also für geeignete Funktionen $f$ und $g$ den Ausdruck
$$
\langle f, g \rangle := \int f \cdot g \d x.
$$
Falls nicht anders vermerkt, bezeichnet $n$ immer eine von Null verschiedene natürliche Zahl.

\section{Glatte Funktionen und der Raum der Testfunktionen}
\label{subsec:smooth}

Ziel dieses Abschnittes ist es, die nötigen Begriffe und Definitionen im Zusammenhang mit glatten Funktionen bereitzustellen.
Im Folgenden bezeichne $\Omega \subseteq \R^n$ stets ein Gebiet, also eine nichtleere, offene und zusammenhängende Teilmenge.

Sei $k \in \N$. Mit $C^k(\Omega)$ bezeichnen wir die Menge aller Funktionen $u$, deren partielle Ableitungen $D^\alpha u$ für alle Multiindices $\alpha \in \N_0^n$ mit $0 \leq |\alpha| \leq k$ existieren und stetig sind. 
Mit
$$
  C^\infty(\Omega) := \bigcap_{k = 0}^\infty C^k(\Omega)
$$
bezeichnen wir den Raum der glatten Funktionen auf $\Omega$.

Im Kontext der Approximation von $L^q$-Funktionen spielt ein Funktionenraum eine wichtige Rolle: der Raum der glatten Funktionen mit kompaktem Träger
$$
C_0^\infty(\Omega) := \{ u \in C^\infty(\Omega) \mid \supp u \text{ kompakt und } \supp u \subseteq \Omega\}.
$$
Wir werden diesem Raum in Abschnitt \ref{subsec:distributionsSobolev} nochmals als Raum der Testfunktionen begegnen.
Zudem werden wir den Raum $C^\infty(\overline\Omega)$ aller Restriktionen $u|_{\overline\Omega}$ von Funktionen aus $C^\infty(\R^n)$ mit
\begin{displaymath}
  \sup_{\alpha  \in \N^n, x \in \R^n} |D^\alpha u(x)| < \infty 
\end{displaymath}
benötigen.
Aufgrund dieser Eigenschaft lässt sich der Raum $C^\infty(\overline\Omega)$ mit einer Norm ausstatten:
\begin{equation}
  \label{eq:cInftyNorm}
  \norm{u}_{C^\infty} := \norm{u}_{C^\infty(\overline\Omega)} := \sup_{|\alpha| \in \N^n, x \in \overline\Omega} |D^\alpha u(x)|.
\end{equation}

Alle eingeführten Räume lassen sich auf natürliche Weise zu Räumen von Vektor\-feldern verallgemeinern.
Man erhält so
\begin{align*}
  C^\infty(\Omega)^m &:= \{(u_1,\dots,u_m) \mid u_j \in C^\infty(\Omega), j = 1,\dots,m\}, \\
  C_0^\infty(\Omega)^m &:= \{(u_1,\dots,u_m) \mid u_j \in C_0^\infty(\Omega), j = 1,\dots,m\} \quad\text{und} \\
  C^\infty(\overline\Omega)^m &:= \{(u_1,\dots,u_m) \mid u_j \in C^\infty(\overline\Omega), j = 1,\dots,m\},
\end{align*}
wobei der letzte Vektorraum durch die Norm
$$
\norm{u}_{C^\infty} := \norm{u}_{C^\infty(\overline\Omega)}^m := \sup_{j = 1,\dots, m} \norm{u_j}_{C^\infty(\overline\Omega)}
$$
zu einem normierten Vektorraum wird.

In der Literatur findet man eine Reihe anderer Definitionen des Symbols $C^\infty(\overline\Omega)$.
Für beschränkte Gebiete stimmen sie jedoch alle überein, wie wir im Folgenden beweisen wollen.
Dazu benötigen wir ein Fortsetzungsresultat von \stein\ \cite[S.172, Proposition 2.2]{stein1970singular}.

\begin{prop}
  \label{prop:stein}
  Sei $f$ eine Funktion auf der abgeschlossenen Menge $F \subseteq \R^n$.
  Dann existiert eine Fortsetzung $\mathcal{E}_0(f)$ von $f$ auf $\R^n$.
  Ist $f$ zusätzlich stetig auf $F$, so ist die Fortsetzung $\mathcal{E}_0$ stetig auf $\R^n$ und glatt auf $F^c$.
\end{prop}

\newpage
\begin{lem}
  \label{lem:CInftyClosedOmega}
  Sei $\Omega \subseteq \R^n$ ein beschränktes Gebiet.
  Für eine Funktion $u \colon \overline\Omega \to \R$ schreiben wir:
  \begin{enumerate}[(a)]
    \item $u \in C^\infty(\overline\Omega)$,
      falls eine Funktion $f \in C^\infty(\R^n)$ mit $\sup_{\alpha \in \N^n, x \in \R^n} | D^\alpha f(x) | <  \infty$ existiert mit $u = f|_{\overline\Omega}$. \cite[S.23, I.3.1]{sohr2001navier}

    \item $u \in C^\infty_{\bullet}(\overline\Omega)$,
      falls eine offene Obermenge $\overline\Omega \subseteq U \subseteq \R^n$ und eine glatte Fortsetzung $\tilde u \colon U \to \R$ von $u$ existieren.

    \item $ u \in C^\infty_{\bullet\bullet}(\overline\Omega)$,
    falls $u|_\Omega \in C^\infty(\Omega)$ gilt und sich für alle Multiindices $\alpha \in \N^n$ die Funktion $D^\alpha u$ stetig auf $\overline\Omega$ fortsetzen lässt. \cite[S.10, 1.28]{adams2003sobolev}, \cite[S.35, II.1.3]{galdi2011navier}
  \end{enumerate}

  Dann gilt
  $$
    C^\infty(\overline\Omega) =
    C^\infty_{\bullet}(\overline\Omega) =
    C^\infty_{\bullet\bullet}(\overline\Omega).
  $$
\end{lem}

\begin{proof}
  %Zunächst ist $\overline\Omega$ als beschränkte abgeschlossene Teilmenge von $\R^n$ kompakt.
  Wir erhalten die Inklusion
  $$
  C^\infty(\overline\Omega) \subseteq C_{\bullet},
  $$ 
  da jedes $u\in C^\infty(\overline\Omega)$ bereits eine glatte Fortsetzung auf die offene Obermenge $\R^n$ besitzt.

  Die Inklusion $$C^\infty_{\bullet}(\overline\Omega) \subseteq C^\infty_{\bullet\bullet}(\overline\Omega)$$ folgt direkt aus der Definition.
  
  Die letzte Inklusion $$C^\infty_{\bullet\bullet}(\overline\Omega) \subseteq C^\infty(\overline\Omega)$$ erfordert mehr Arbeit.
  Sei dazu $u \in C^\infty_{\bullet\bullet}(\overline\Omega)$.
  %Es ist eine Eigenschaft von \hausdorff\hyp{}Räumen, dass stetige Abbildungen höchstens eine Fortsetzung auf den Abschluss ihres Definitionsberechs besitzen \cite[S.112, Lemma 4.2.4]{bartsch2015topologie}.
  %Wir wollen daher $u$ mit seiner Fortsetzung auf $\overline\Omega$ identifizieren.
  Das Fortsetzungsresultat von \stein\ \ref{prop:stein} erlaubt es, die stetige Funktion $u$ ausgehend von der abgeschlossenen Menge $\overline\Omega$ auf ganz $\R^n$ zu einer Funktion $\tilde u$ fortzusetzen.
  Diese Fortsetzung besitzt jedoch im Allgemeinen keine beschränkten Ableitungen wie in der Definition von $C^\infty(\overline\Omega)$ gefordert.
  Wir werden daher $\tilde u$ außerhalb von $\overline\Omega$ geeignet modifizieren.
  Da $\Omega$ als beschränkt vorausgesetzt wurde, existiert ein offener Ball $B_r(0)$ mit Radius $r > 0$ und 
  $$
   \Omega \subset\subset B_r(0).
  $$ 
  Wir multiplizieren nun die Fortsetzung $\tilde u$ mit einer glatten \emph{cut-off}\hyp{}Funktion $\psi$, für die
  $$ 
  \psi(x) = \begin{cases}
                          1 &\text{ für alle } x \in \overline\Omega \\
                          0 &\text{ für alle } x \in \R^n \setminus B_r(0)
             \end{cases}
  $$
  und folglich $\supp \psi \subseteq B_r(0)$ gilt.
  Eine solche Funktion erhält man zum Beispiel aus der Faltung der charakteristischen Funktion $\chi_{\overline\Omega}$ mit einem Glättungskern, wie der Beweis des \urysohn\hyp{}Lemmas für glatte Funktionen in \cite[S.88, Proposition 6.5]{folland2009ra} zeigt. 
  Für $f := \tilde u \cdot \psi$ erhalten wir
  $$
  \sup_{|\alpha| < \infty, x \in \R^n} | D^\alpha f |
   = \sup_{|\alpha| < \infty, x \in \overline B_r} | D^\alpha f |
   < \infty,
  $$
  da $\overline B_r$ kompakt ist.
  Aufgrund der Fortsetzungseigenschaft gilt weiterhin $f|_{\overline\Omega} = u$ und damit ist $f$ die für (a) gesuchte Fortsetzung.
\end{proof}

Zuletzt erwähnen wir den Untervektorraum der divergenzfreien glatten Vektorfelder
$$
C_{0,\sigma}^\infty(\Omega) := \{u \in C_0^\infty(\Omega)^n \mid \div u = 0\},
$$
 welcher in der schwachen Lösungstheorie der stationären inkompressiblen \navier\hyp{}\stokes\hyp{}Gleichungen als Testraum eingesetzt wird.

\section{Schwache Differenzierbarkeit}
\label{subsec:distributionsSobolev}

Wir wenden uns nun einer Verallgemeinerung des Differenzierbarkeitsbegriffs zu, welcher für die Lösungstheorie partieller Differentialgleichungen von zentraler Bedeutung ist: die \emph{schwache} Differenzierbarkeit.
Es bezeichne im Folgenden wieder $\Omega \subseteq \R^n$ ein Gebiet mit $n \geq 1$.

\subsection{Lokal integrierbare Funktionen}

Wir werden im Folgenden oft eine Obermenge der bezüglich des \lebesgue\hyp{}Maßes auf $\R^n$ integrierbaren Funktionen verwenden:
Für $1 \leq q \leq \infty$ schreiben wir
$$
  u \in L^q_{\loc}(\Omega)
$$
und nennen $u$ \emph{lokal integrierbar}, falls $u \in L^q(B)$ für alle offenen Bälle $B \subseteq \Omega$ mit $\overline B \subseteq \Omega$ gilt.

\begin{bem}
  Eine Funktion $u$ ist genau dann lokal integrierbar über $\Omega$, falls $u \in L^q(K)$ für alle Kompakta $K \subseteq \Omega$ gilt.

  Diese Aussage findet man ebenfalls in der Literatur zur Definition lokaler Integrierbarkeit.
  Tatsächlich ist sie äquivalent zur obigen Definition.
  Denn einerseits ist für alle Bälle $B$ auch $\overline B$ ein Kompaktum und die Aussage folgt aus der Inklusionsbeziehung $L^q(\overline B) \subseteq L^q(B)$.
  Andererseits lässt sich jedes Kompaktum $K$ durch endlich viele Bälle $B_1, \dots, B_n$ mit $\overline B_i \subseteq \Omega, i =1,\dots,m,$ überdecken. 
  Gilt nun $u \in L^q(B_i)$ für alle $i = 1, \dots, m$, so gilt insbesondere $u \in L^q(\bigcup_{i = 1}^m B_i)$ und die Umkehrung der Aussage folgt aus der Inklusionsbeziehung $L^q(\bigcup_{i = 1}^m B_i) \subseteq L^q(K)$.
\end{bem}

Des Weiteren schreiben wir
$$
u \in L^q_{\loc}(\overline\Omega),
$$
falls $u \in L^q(B \cap \Omega)$ für jeden offenen Ball $B \subseteq \R^n$ mit $B \cap \Omega \neq \emptyset$ gilt.
Zusammenfassend gilt also die folgende Inklusionsbeziehung:
$$
L^q(\Omega) \subseteq L^q_{\loc}(\overline\Omega) \subseteq L^q_{\loc}(\Omega).
$$

\subsection{Distributionen}

In Abschnitt \ref{subsec:smooth} hatten wir bereits den Raum der Testfunktionen $C_0^\infty(\Omega)$ kennengelernt.
Wir werden uns vor allem für seinen Dualraum, den Raum der stetigen Funktionale auf $C_0^\infty(\Omega)$, interessieren.
Um überhaupt über Stetigkeit von Funktionalen auf $C_0^\infty(\Omega)$ sprechen zu können, müssen wir jedoch zunächst eine Topologie auf dem besagten Raum festlegen.
Dabei beschränken wir uns auf den aus dieser Topologie resultierenden Stetigkeitsbegriff:
Ein lineares Funktional $F \colon C_0^\infty(\Omega) \to \R$ ist stetig, genau dann, wenn für jedes beschränkte Teilgebiet $G \subseteq \Omega$ mit $\overline G \subseteq \Omega$ ein $k \in \N_0$ und $C = C(F,G)$ existiert, sodass
$$
|F(\varphi)| \leq C \norm{\varphi}_{C^k(\overline G)}
$$
gilt, wobei $\norm{\cdot}_{C^k(\overline G)}$ gerade die in Gleichung (\ref{eq:cInftyNorm}) definierte Norm bezeichnet.
Die hinter diesem Stetigkeitsbegriff stehende Topologie macht $C_0^\infty(\Omega)$ zu einem lokalkonvexen Vektorraum.
Eine ausführlichere Beschreibung der beschriebenen Topologie findet sich in \cite[S.433f.]{werner2011fa}.
Den Raum $C_0^\infty(\Omega)'$ aller im obigen Sinne stetigen Funktionale
$$
F \colon C_0^\infty(\Omega) \to \R, \quad \varphi \mapsto F(\varphi) =: [F, \varphi],
$$
bezeichnen wir als den Raum der \emph{Distributionen}.
Hierbei stellt $[\cdot,\cdot]$ die duale Paarung auf $\Omega$ dar. 

Wir betrachten nun spezielle Distributionen.
Ist $f \in L^1_{\loc}$, so verstehen wir unter 
$$
f \mapsto \langle f, \varphi \rangle := \int_\Omega f \varphi \d x, \quad \varphi \in C_0^\infty(\Omega), 
$$
die zu $f$ gehörige \emph{reguläre} Distribution.
Dass diese Zuordnung sogar injektiv ist und somit zur Inklusionsbeziehung $L^1_{\loc}(\Omega) \subseteq C^\infty_0(\Omega)'$ führt, lässt sich in \cite[S.432, Beispiel (a)]{werner2011fa} nachlesen.

\subsection{Differentiation von Distributionen}

Distributionen ermöglichen es, den für die Analysis grundlegenden Begriff der Differenzierbarkeit zu verallgemeinern.
Ist $\alpha \in \N_0^n$ ein Multiindex und $D^\alpha \colon C_0^\infty(\Omega) \to C_0^\infty(\Omega)$ der zugehörige Ableitungsoperator, so definieren wir für eine Distribution $F \in C_0^\infty(\Omega)'$ ihre distributionelle Ableitung $D^\alpha F \in C_0^\infty(\Omega)'$ durch
$$
  [D^\alpha F, \varphi] := (-1)^{|\alpha|} [F, D^\alpha \varphi].
$$
Bis auf ein Vorzeichen stimmt also die distributionelle Ableitung mit der zum Ableitungsoperator dualen Abbildung $(D^\alpha)' \colon C_0^\infty(\Omega)' \to C_0^\infty(\Omega)'$ überein.
Dass die distributionelle Ableitung wohldefiniert und sogar schwach*-stetig ist, lässt sich in \cite[S.434, Lemma VIII.5.7]{werner2011fa} nachlesen.

Im Unterabschnitt zu glatten Funktionen hatten wir die dortigen Definitionen auf Produkträume ausgeweitet.
Wir versehen den so entstandenen Produktraum $C_0^\infty(\Omega)^m$ mit der Produkttopologie und wenden uns nun dem zugehörigen Distributionenraum $(C_0^\infty(\Omega)^m)'$ zu.

Wir betrachten dazu die Distribution
$$
  F = (F_1, \dots, F_m), \quad F_j \in C_0^\infty(\Omega)', \quad j = 1, \dots, m,
$$
und definieren für alle 
$$\varphi = (\varphi_1, \dots, \varphi_m) \in C_0^\infty(\Omega)^m$$
die Distribution
\begin{equation}
  \label{eq:pairingVector}
  [F, \varphi] := [F_1, \varphi_1] + \dots + [F_m, \varphi_m].
\end{equation}
Durch Betrachtung von Funktionen des Typs 
$$\psi^{(i)} = (0, \dots,0, \psi_i, 0,\dots,0) \in C_0^\infty(\Omega)^m, \quad i \in \{1,\dots,m\}, $$
erkennen wir, dass jedes bezüglich der Produkttopologie stetige Funktional auf $C_0^\infty(\Omega)^m$ diese Form besitzt.
Für den Raum der Distributionen auf den Testfunktionen $C_0^\infty(\Omega)^m$ gilt also
\begin{align*}
  (C_0^\infty(\Omega)')^m 
  &= (C_0^\infty(\Omega)^m)' \\
  &= \{(F_1,\dots,F_m) \mid F_j \in C_0^\infty(\Omega)', j = 1, \dots,m\}.
\end{align*}

Wir erinnern uns an den im Unterabschnitt zu glatten Funktionen definierten Funktionenraum $C_{0, \sigma}^\infty(\Omega) \subseteq C_0^\infty(\Omega)^n$ der divergenzfreien Funktionen.
Nach dem Fortsetzungssatz von \hahn\hyp{}\banach\ für lokalkonvexe Vektorräume (\cite[S.408, Satz VIII.2.8]{werner2011fa}) erhalten wir die stetigen Funktionale auf $C_{0,\sigma}^\infty$ gerade als Einschränkungen der stetigen Funktionale auf $C_0^\infty(\Omega)^n$. 
Es gilt somit
$$
C_{0,\sigma}^\infty(\Omega)' = \{F|_{C_{0,\sigma}^\infty(\Omega)} \mid F \in (C_0^\infty(\Omega)^n)'\}.
$$

Betrachten wir nun den \hilbert\hyp{}Raum $L^2(\Omega)^n$.
Wie Lemma \ref{lem:mollification} noch zeigen wird, liegen die glatten Funktionen mit kompaktem Träger dicht in $L^2(\Omega)$. 
Ebenso stellt der Produktraum $C_0^\infty(\Omega)^n$ einen dichten Unterraum von $L^2(\Omega)^n$ dar.
Um diese Tatsache zu imitieren, definieren wir den Unterraum 
$$
  L^2_\sigma(\Omega) 
  := \overline{C_{0,\sigma}^\infty(\Omega)}^{\norm{\cdot}_2}
  \subseteq L^2(\Omega)^n.
$$
Identifizieren wir nun jede Funktion $u \in L^2(\Omega)^n$ mit dem Funktional
$$
\langle u, \cdot \rangle \colon \varphi \mapsto \langle u, \varphi \rangle, \quad \varphi \in C_0^\infty(\Omega)^n,
$$
so erhalten wir die Einbettung 
$$
  L^2(\Omega)^n \subseteq (C_0^\infty(\Omega)^n)'.
$$
Ebenso lassen sich Funktionen $u \in L^2_\sigma(\Omega)$ mit der entsprechenden Einschränkung
$$
\langle u, \cdot \rangle \colon \varphi \mapsto \langle u, \varphi \rangle, \quad \varphi \in C_{0,\sigma}^\infty(\Omega)^n
$$
identifizieren, was die Einbettung
$$
  L^2_\sigma(\Omega) \subseteq (C_{0,\sigma}^\infty(\Omega)^n)'
$$
liefert.

\subsection{Sobolev-Räume}

Im Allgemeinen wird die Ableitung einer regulären Distribution nicht wieder regulär sein.
Wir definieren daher für alle $k \in \N$ und $1 \leq q \leq \infty$ den $L^q$\hyp{}\sobolev\hyp{}Raum $W^{k,q}(\Omega)$ der Ordnung $k$ durch
$$
W^{k,q}(\Omega) := \{ u \in L^q(\Omega) \mid D^\alpha u \in L^q(\Omega) \text{ für alle } |\alpha| \leq k\}.
$$
Für $u \in W^{k,q}(\Omega)$ bezeichnen wir $D^\alpha u \in L^q(\Omega)$ als \emph{schwache Ableitung} von $u$.
Hierbei identifizieren wir die reguläre Distribution $D^\alpha u$ immer direkt mit der korrespondierenden $L^q$\hyp{}Funktion.

Wir machen $W^{k,q}$ zu einem normierten Vektorraum durch die \sobolev\hyp{}Norm
$$
  \norm{u}_{W^{k,q}(\Omega)} 
  := \norm{u}_{W^{k,q}} 
  := \norm{u}_{k,q} 
  := \begin{cases} 
    \left(\sum_{|\alpha| \leq k} \norm{D^\alpha u}_q^q \right)^\frac{1}{q}\quad&\text{für } 1 \leq q < \infty \\
    \max_{|\alpha| \leq k} \norm{D^\alpha u}_\infty \quad&\text{für }q = \infty.
    \end{cases}
$$

\begin{lem}
  Für ein Gebiet $\Omega \subseteq \R^n$ mit $n \geq 1$ und $1 \leq q \leq \infty$ ist $W^{k,q}(\Omega)$ ausgestattet mit der \sobolev\hyp{}Norm ein \banach\hyp{}Raum.
\end{lem}

\begin{proof}
  Sei $(u_j)_{j \in \N}$ eine \cauchy\hyp{}Folge in $W^{k,q}(\Omega)$.
  Nach Definition ist dann auch $(D^\alpha u_j)_{j \in \N}$ für alle $|\alpha| \leq k$ eine \cauchy\hyp{}Folge in $L^q$.
  Folglich existieren für alle $|\alpha| \leq k$ Funktionen $u_\alpha \in L^q$ mit 
  $$
  \lim_{j \to \infty} \norm{D^\alpha u_j - u_\alpha}_q = 0.
  $$
  Sei $u_0 := u_{(0,\dots,0)}$.
  Wir behaupten nun, dass im distributionellen Sinne die Identität $u_\alpha = D^\alpha u_0$ für alle $\alpha \in \N^n$ gilt.
  Wir zeigen dazu, dass für alle $\varphi \in C_0^\infty(\Omega)$ die Identität
  $$
  [u_\alpha, \varphi] = [D^\alpha u_0, \varphi]
  $$
  gilt.
  Dazu halten wir zunächst fest, dass mit der \hoelder\hyp{}Ungleichung
  $$
  \int_\Omega |(D^\alpha u_j - u_\alpha)| \varphi \d x
  \leq \norm{(D^\alpha u_j - u_\alpha)}_q \norm{\varphi}_{q'} \to 0, \quad j \to \infty
  $$
  sowie
  $$
  \int_\Omega |(u_j - u_0) D^\alpha \varphi| \d x
  \leq \norm{(u_j - u_0)}_q \norm{D^\alpha \varphi}_{q'} \to 0, \quad j \to \infty
  $$
  folgen, wobei $q'$ den zu $q$ konjugierten Exponenten bezeichne.
  Hieraus folgen die Identitäten
  $$
  \lim_{j \to \infty} [D^\alpha u_j, \varphi] = \lim_{j \to \infty} \int_\Omega D^\alpha u_j \varphi \d x =  \int_\Omega u_\alpha \varphi \d x =  [u_\alpha, \varphi] 
  $$
  und
  $$
    \lim_{j \to \infty} [u_j, D^\alpha\varphi] = \lim_{j \to \infty} \int_\Omega u_j D^\alpha\varphi \d x =  \int_\Omega u_0 D^\alpha \varphi \d x = [u_0, D^\alpha \varphi].
  $$
  Dies impliziert schließlich
  $$ 
  [u_\alpha, \varphi] 
  = \lim_{j \to \infty} [D^\alpha u_j, \varphi] \\
  = \lim_{j \to \infty} (-1)^{|\alpha|}[u_j, D^\alpha \varphi] \\
  =  (-1)^{|\alpha|}[u_0, D^\alpha \varphi]
  = [D^\alpha u_0, \varphi].
  $$
  Daraus folgt jedoch gerade die Behauptung $D^\alpha u_0 = u_\alpha$.
  Insgesamt ergibt sich also $\lim_{j \to \infty} \norm{u_j - u_0}_{k,q} =  0$, was zu beweisen war.
\end{proof}

Wir definieren den Unterraum
$$
  W_0^{k,q}(\Omega) := \overline{C_0^\infty(\Omega)}^{\norm{\cdot}_{k,q}},
$$
welcher die glatten Funktionen mit kompaktem Träger als dichte Teilmenge bezüglich der \sobolev\hyp{}Norm enthält.

Bisher hatten wir nur \sobolev\hyp{}Räume positiver Ordnung betrachtet.
Wir definieren für $1 < q < \infty$ die \sobolev\hyp{}Räume negativer Ordnung 
$$
  W^{-k,q}(\Omega) := W_0^{k,q'}(\Omega)'.
$$
Wie für den Dualraum eines normierten Raumes üblich, ist $W^{-k,q}(\Omega)$ mit Operatornorm 
$$
\norm{F}_{W^{-k,q}(\Omega)} := \norm{F}_{W^{-k,q}} := \norm{F}_{-k,q} := \sup_{0 \neq \varphi \in C_0^\infty(\Omega)} \frac{|[F,\varphi]|}{\norm{\varphi}_{k,q'}}
$$
ein \banach\hyp{}Raum.
Wir beachten, dass die Operatornorm bereits durch die bezüglich \sobolev\hyp{}Norm in $W^{k,q}(\Omega)$ dichte Teilmenge $C_0^\infty(\Omega)$ eindeutig festgelegt ist.

Es ist möglich einen \sobolev\hyp{}Raum $W^{k,q}(\Omega)$ für $1 < q < \infty$ als abgeschlossenen Teilraum eines $L^q$\hyp{}Raumes zu realisieren \cite[S.61, 3.5]{adams2003sobolev}.
Als solcher ist er insbesondere reflexiv.
Wir halten diese wichtige Eigenschaft in einem Lemma fest.

\begin{lem}
  Für ein Gebiet $\Omega \subseteq \R^n$ mit $n \geq 1$ und $1 < q < \infty$ ist $W^{k,q} = W^{k,q}(\Omega)$ ausgestattet mit der \sobolev\hyp{}Norm ein reflexiver Raum.
\end{lem}

Aus der Reflexivität der \sobolev\hyp{}Räume leiten wir für $1 < q < \infty$ die folgende Identität ab:
$$
  W_0^{k,q'}(\Omega) = W^{-k,q}(\Omega)'.
$$
Wir können also jede Funktion $u \in W_0^{k,q'}(\Omega)$ mit dem Funktional
$$
[\cdot, u] \colon F \mapsto [F, u], \quad F \in W^{-1,q}(\Omega)
$$
identifizieren.

%Wir definieren nun die $W^{k,q}_{\loc}$\hyp{}Räume.
Wir bezeichnen mit $W^{k,q}_{\loc}(\Omega)$ den Raum aller Funktionen $u$ mit $D^\alpha u \in L^q_{\loc}(\Omega)$ für alle Multiindices $\alpha \in \N_0^n$ mit $|\alpha| \leq k$.
Des Weiteren definieren wir $W^{k,q}(\overline\Omega)$ als den Raum aller Funktionen $u$ mit $D^\alpha u \in L^q_{\loc}(\overline\Omega)$.
Nun zu den \sobolev\hyp{}Räumen negativer Ordnung. 
Wir definieren den Raum $W^{-k,q}_{\loc}(\Omega)$ als den Raum aller Distributionen
$$
F \colon \varphi \mapsto [F,\varphi], \quad \varphi \in C_0^\infty(\Omega),
$$
sodass für die Operatornorm die Ungleichung
\begin{equation}
  \label{eq:operatornorm}
\norm{F}_{-k,q} := \sup_{0 \neq \varphi \in C_0^\infty(\Omega_0)}  \frac{|[F,\varphi]|}{\norm{\varphi}_{k,q'}} < \infty
\end{equation}
für alle beschränkten Teilgebiete $\Omega_0 \subseteq \Omega$ mit $\overline\Omega_0 \subseteq \Omega$ gilt.

Wie schon im Unterabschnitt zu glatten Funktionen vorgestellt, können wir auch die Definition der \sobolev\hyp{}Räume auf Produkträume ausdehnen. 
Wir definieren hierzu
$$
W^{k,q}(\Omega)^m := \{ (u_1,\dots,u_m) \mid u_j \in W^{k,q}(\Omega), \quad j = 1,\dots, m\}
$$
und versehen den so entstandenen Vektorraum mit der Norm
$$
\norm{u}_{W^{k,q}(\Omega)^m} 
:= \norm{u}_{k,q}
:= \left( \sum_{j = 1}^m \norm{u_j}_{k,q}^q \right)^{\frac{1}{q}},
$$
wodurch dieser zu einem \banach\hyp{}Raum wird.


%\newpage
\section{Eigenschaften der Faltung}
\label{subsec:mollification}

Aus der Integrationstheorie ist bekannt, dass sich $L^q$-Funktionen durch Faltungen mit Glättungskernen $\mathcal{F}_\varepsilon$ approximieren lassen.
Für unsere Zwecke wird es dafür ausreichend sein den Standardkern
$$
\mathcal{F}(x) := \begin{cases}
                     c \cdot \exp\left(-\frac{1}{1 - \norm{x}}\right), \quad &\norm{x} < 1 \\
                     0, \quad  &\text{sonst}
                  \end{cases} \quad,\quad x \in \R^n
$$
zu betrachten und
$$
\mathcal{F}_\varepsilon(x) := \frac{1}{\varepsilon^n} \mathcal{F}\left(\frac{x}{\varepsilon}\right)
$$ 
zu setzen.
Hierbei sei $c$ so gewählt, dass $\int_{B_1(0)} \mathcal{F} \d x = 1$ gilt.

Folgendes Lemma fasst nochmals die Approximationseigenschaft von Glättungskernen zusammen. Ein Beweis findet sich in \cite[S.36, Theorem 2.29(c)]{adams2003sobolev}.
\begin{lem}
  \label{lem:mollification}
  Sei $\Omega \subseteq \R^n$ mit $n \geq 1$ ein Gebiet und $1 \leq q < \infty$ sowie $\varepsilon > 0$.
  Für alle $u \in L^q(\Omega)$ gilt dann
  $$
  \norm{(\mathcal{F}_\varepsilon \ast u)}_{L^q(\Omega)} \leq  \norm{u}_{L^q(\Omega)}
  $$
  und
  \begin{displaymath}
    \lim_{\varepsilon \to 0} \norm{(\mathcal{F}_\varepsilon \ast u) - u}_{L^q(\Omega)} = 0. 
  \end{displaymath}
\end{lem}

Sei nun $\Omega_0 \subseteq \overline\Omega_0 \subseteq \Omega$ ein beschränktes Teilgebiet und
\begin{equation}
  \label{eq:convolutionEpsilon}
  0 < \varepsilon < \dist(\Omega_0, \partial\Omega).
\end{equation}
Zusätzlich sei $u \in L^1_{\loc}(\Omega)$.
Setzt man $u(x) := 0$ für alle $x \in \R^n \setminus \Omega$, so folgt $u \in L^1_{\loc}(\R^n)$.
Wie \cite[S.171, Theorem 6.30(b)]{rudin1991fa} zeigt, lassen sich aus der Integrationstheorie bekannte Eigenschaften der Glättungen auch auf Faltungen 
$$
u^\varepsilon := \mathcal{F}_\varepsilon \ast u
$$
von Glättungskernen $\mathcal{F}_\varepsilon \ast u$ mit Distributionen $u \in C_0^\infty(\Omega)'$ verallgemeinern.
Einerseits zählt dazu, dass $u^\varepsilon \in C^\infty(\R^n)$ gilt.
Da $\Omega_0$ als beschränkt vorausgesetzt wurde, gilt nach Lemma \ref{lem:CInftyClosedOmega}(b$\Rightarrow$a) sogar $u^\varepsilon \in C^\infty(\overline\Omega_0)$.
Andererseits ist für alle Multiindices $\alpha \in \N_0^n$ und alle $x \in \Omega_0$ die folgende Gleichung gültig:
\begin{equation}
  \label{eq:convolutionDiff}
  (D^\alpha u^\varepsilon)(x) 
  = (\mathcal{F}_\varepsilon \ast (D^\alpha u))(x)
  = ( (D^\alpha \mathcal{F}_\varepsilon) \ast u)(x), \quad x \in \Omega_0.
\end{equation}

\begin{lem}
\label{lem:nablauconstant}
  Sei $\Omega \subseteq \R^n$ mit $n \geq 1$ ein Gebiet und $1 \leq q < \infty$.
  Ist $u \in L^1_{\loc}(\Omega)$ und gilt $\nabla u = 0$ im distributionellen Sinne, dann ist $u$ konstant.
\end{lem}

\begin{proof}
  Für alle beschränkten Teilgebiete $\Omega_0 \subseteq \Omega$ und $\varepsilon$ wie in Ungleichung (\ref{eq:convolutionEpsilon}) folgern wir unter Verwendung der Gleichung (\ref{eq:convolutionDiff})
  \begin{align*}
    \nabla u^\varepsilon
    &= (D_1 (\mathcal{F}_\varepsilon \ast u), \dots, D_n (\mathcal{F}_\varepsilon \ast u))^T \\
    &= ( \mathcal{F}_\varepsilon \ast D_1 u, \dots, \mathcal{F}_\varepsilon \ast D_n u)^T \\
    &= (0, \dots, 0)^T.
  \end{align*}
  Es ist $u^\varepsilon$ eine glatte Funktion.
  Zudem ist $\Omega_0$ als offene zusammenhängende Teilmenge von $\R^n$ insbesondere wegzusammenhängend.
  Ausgehend von der Integraldarstellung des Funktionszuwachses \cite[S.57, (11)]{koenigsberger2004ana2} gilt $u_\varepsilon = C_\varepsilon$ auf ganz $\Omega_0$.
  Mit Lemma \ref{lem:mollification} folgt nun die Netzkonvergenz $C_\varepsilon \to C$ für $\varepsilon \to 0$ auf $\Omega_0$.
  Wir greifen auf eine im Abschnitt \ref{sec:lipschitzDomains} zu \lipschitz\hyp{}Gebieten eingeführte Konstruktion vor:
  Aufgrund von Lemma \ref{lem:lipschitzExhaust} lässt sich das Gebiet $\Omega$ von einer Folge $(\Omega_j)_{j \in \N}$ offener beschränkter \lipschitz\hyp{}Gebiete ausschöpfen. 
  Durch Anwendung des für das Gebiet $\Omega_0$ beschriebenen Arguments auf die Folgenglieder $\Omega_j$ erhalten wir somit $u = C$ auf ganz $\Omega$. 
  Hieraus folgt die Behauptung.
\end{proof}



\chapter{Lösungen von $\nabla p = f$}
\section{Lipschitz-Gebiete und Gebietsapproximation}

Im Folgenden sei $\R^n$ immer mit der \euklid ischen Metrik versehen. Es bezeichne zudem
$$
\dist(X,Y) := \inf\{ \norm{x - y} \mid x \in X, y \in Y\}
$$
den \euklid ischen Abstand zweier Mengen $X,Y \subseteq \R^n$.

\begin{defn}
  \lipschitz\hyp{}Gebiete  
\end{defn}

\begin{lem}
  \label{lem:distBoundary}
  Seien $\emptyset \subsetneq A,B \subsetneq \R^n$. Gilt $A \subseteq B$, so folgt $\dist(a,\partial A) \leq \dist(a, \partial B)$ für alle $a \in A$.
\end{lem}

\begin{proof}
  Da $\R^n$ ein zusammenhängender Raum ist, sind die einzigen Mengen mit leerem Rand die leere Menge und der ganze Raum.
  Dies wurde jedoch in der Voraussetzung des Lemmas bereits ausgeschlossen, daher nimmt $\dist$ nur endliche Werte an.

  Sei nun $a \in A$ und $b \in \partial B$
  Wir betrachten den Strahl $s \colon [0,1] \to \mathbb{R}^n$ mit $s(0) = a$ und $s(1) = b$.
  Zudem können wir annehmen, dass $a,b \not\in \partial A$ gilt, da ansonsten die Ungleichung sofort erfüllt ist.
  Wir wollen im Folgenden nachweisen, dass ein $t' \in (0,1)$ mit $s(t') \in \partial A$ existiert.

  Dazu definieren wir die Funktion
  $$
  f(x) := (1 - 2\chi_A(x)) \cdot \dist(x, \partial A), \quad x \in \R^n,
  $$
  wobei $\chi_A$ die charakteristische Funktion der Menge $A$ bezeichne.
  Als Nächstes weisen wir nach, dass $f$ auf $\R^n$ stetig ist.
  Es ist bekannt, dass die Funktion $\dist(\cdot, \partial A)$ stetig ist \cite{koenigsberger}.
  Auf $\R^n \setminus \partial A$ ist zudem $(1 - 2\chi_A(x))$ konstant gleich $-1$ beziehungsweise $1$, also ist $f$ dort stetig.

  Sei nun $(x_n)_{n \in \N}$ eine Folge in $\R^n$ mit $\lim_{n \to \infty} x_n = x \in \partial A$.
  Dann gilt
  $$
  |(1 - 2\chi_A(x_n)) \cdot \dist(x_n, \partial A)| 
  \leq \dist(x_n,\partial A)
  \to \dist(x,\partial A) 
  = 0,
  $$
  also auch 
  $$
  \lim_{n \to \infty} f(x_n) = 0 = f(x),
  $$
  da aufgrund der Abgeschlossenheit von $\partial A$ die Gleichheit $\dist(x,\partial A) = 0$ äquivalent zu $x \in \partial A$ ist ($\ast$).

  Auf $f \circ s$ lässt sich nun der Zwischenwertsatz anwenden, denn nach Voraussetzung gelten
  $$
  (f \circ s)(0) = f(a) = -1 \quad\text{und}\quad (f \circ s)(1) = f(b) = 1.
  $$
  Somit existiert also ein $t' \in (0,1)$ mit $(f\circ s)(t') = f(s(t')) = 0$, was nach ($\ast$) äquivalent zu $s(t') \in \partial A$ ist.
  Es ist also
  $$
  \norm{a - b} \geq \norm{a - s(t')} - \norm{s(t') - b} \geq \norm{a - s(t')} \geq \dist(x,\partial A).
  $$
  Da dies für alle $b \in \partial B$ gilt folgt sogleich $ \dist(a, \partial B) \geq \dist(a, \partial A)$.
\end{proof}

\begin{bem}
  Wie aus dem Beweis von Lemma \ref{lem:distBoundary} ersichtlich ist, reicht es bereits aus $A$ und $B$ als Teilmengen der konvexen Hülle von $B$ mit euklidischer Spurmetrik und entsprechender Abstandsfunktion $\dist$ zu betrachten.
  Die Konvexität ist hierbei notwendig, wie man an einem Beispiel zeigen kann.
\end{bem}

\begin{lem}
  \label{lem:lipschitzExhaust}
  Sei $\Omega \subseteq \R^n$ mit $n \geq 2$ ein Gebiet.
  Dann existiert eine Folge $(\Omega_j)_{j \in \N}$ beschränkter \lipschitz\hyp{}Gebiete $\Omega_j \subseteq \Omega$ und eine Folge $(\varepsilon_j)_{j \in \N}$ positiver reeller Zahlen mit folgenden Eigenschaften:
  \begin{enumerate}[a)]
    \item Für alle $j \in \N$ gilt $\overline{\Omega}_j \subseteq \Omega_{j + 1}$.
    \item Für alle $j \in \N$ gilt $\varepsilon_{j + 1} \leq \dist(\Omega_j, \partial \Omega_{j + 1}).$
    \item Es gilt $\lim_{j \to \infty} \varepsilon_j = 0$.
    \item Die Gebiete $\Omega_j$ schöpfen $\Omega$ aus.
  \end{enumerate}
\end{lem}

\begin{proof}
Im Folgenden bezeichne $B_r(x) \subseteq \R^n$ den bezüglich \euklid ischer Topologie offenen Ball mit Radius $r$ und Mittelpunkt $x$.

  Für ein festgewähltes $x_0 \in \Omega$ betrachten wir den Schnitt 
  $$
  \Omega' := \Omega \cap B_1(x_0).
  $$ 
  Als Schnitt offener Mengen ist $\Omega'$ wiederum offen. 
  Bezüglich der Teilraumtopologie muss $\Omega'$ jedoch nicht zwingend zusammenhängend sein.
  Wir bezeichnen nun mit $\widetilde\Omega_1$ die Zusammenhangskomponente von $\Omega'$, welche $x_0$ enthält.
  Da die Zusammenhangskomponenten eines topologischen Raumes immer eine Partition desselben bilden, ist $\Omega'$ eindeutig bestimmt.
  Insbesondere gilt für den Rand
  $$ 
  \partial \widetilde\Omega_1 \subseteq \overline{B_1(x_0)}, 
  $$
  er ist somit als abgeschlossene Teilmenge des Kompaktums $\overline{B_1(x_0)}$ selbst kompakt.
  Für alle $\varepsilon > 0$ lässt sich daher $\partial \widetilde\Omega_1$ durch endlich viele Bälle $B_\varepsilon(x_j)$, mit $x_j \in \partial \widetilde\Omega_1$ für alle $j = 1,\dots,m$, überdecken:
  $$ 
  \partial \widetilde\Omega_1 \subseteq \bigcup_{j = 1}^m B_\varepsilon(x_j).
  $$

  Wir definieren nun 
  $$
  \widehat\Omega_1 := \widetilde\Omega_1 \setminus \bigcap_{j = 1}^m \overline{B_\varepsilon(x_j)}
  $$
  und wählen $0 < \varepsilon < 1$ so klein, dass zusätzlich $x_0 \in \widehat\Omega_1$ gilt. 
  Dies lässt sich immer erreichen, da $\widetilde\Omega_1$ als bezüglich Teilraumtopologie offen-abgeschlossene Menge in $\Omega'$ auch in $\R^n$ offen ist und daher ein $\delta > 0$ mit $B_\delta(x_0) \subseteq \Omega'$ existiert.
  Hiermit besitzt bereits ein $\varepsilon < \dist(x_0, \partial\widetilde\Omega_1) - \delta$ die geforderte Eigenschaft.

  Man erkennt nun $\widehat\Omega_1$ als beschränktes \lipschitz\hyp{}Gebiet, da $\partial\widehat\Omega_1$ sämtlich aus Teilen der Ränder der Bälle $B_\varepsilon(x_j)$ besteht.
  Wir setzen nun $\Omega_1 := \widehat\Omega_1$ und $\epsilon_1 := \epsilon$ und führen diese Konstruktion weiter fort.

  Wir wählen wieder 
  $$
  \widetilde\Omega_2 \subseteq \Omega \cap B_2(x_0)
  $$
  als die $x_0$ enthaltende Zusammenhangskomponente des Schnitts von $\Omega$ und $B_2(x_0)$ und konstruieren analog zum ersten Schritt ein Gebiet $\widehat\Omega_2$ mit $0 < \varepsilon < \tfrac{1}{2}$ und $\varepsilon < \dist(\Omega_1, \partial\widetilde\Omega_2)$.
  Dies ist jedoch nur möglich, falls $0 < \dist(\Omega_1, \partial\widetilde\Omega_2)$ gilt, was wir im Folgenden beweisen werden.

  Zunächst gilt nach Konstruktion die Inklusionskette
  $$
  \widehat\Omega_1 \subseteq \widetilde\Omega_1 \subseteq \widetilde\Omega_2.
  $$
  Hieraus folgt mit Lemma \ref{lem:distBoundary}, dass
  \begin{displaymath}
    \dist(x,\partial\widehat\Omega_1) 
    \leq \dist(x,\partial\widetilde\Omega_1) 
    \leq \dist(x,\partial\widetilde\Omega_2) \tag{$\ast$}
  \end{displaymath}
  für alle $x \in \widehat\Omega_1$ gilt.
  Des Weiteren gilt 
  \begin{displaymath}
  0 < \lambda \leq \dist(\widehat\Omega_1,\partial\widetilde\Omega_1) \tag{$\ast\ast$},
  \end{displaymath}
  wobei $\lambda$ die \lebesgue\hyp{}Zahl der Überdeckung $B_{\varepsilon_1}(x_1),\dots,B_{\varepsilon_1}(x_m)$ von $\partial\widetilde\Omega_1$ aus dem ersten Schritt des Beweises bezeichne.
  Die Ungleichungen ($\ast$) und ($\ast\ast$) zusammen ergeben nun die Behauptung.

  Setzen wir noch $\Omega_2 := \widehat\Omega_2$ und $\varepsilon_2 := \varepsilon$, so erhalten wir einerseits $\overline\Omega_1 \subseteq \Omega_2$, denn $\Omega_1 \subseteq \Omega_2$ gilt nach Konstruktion, sowie $0 < \dist(x,\partial\widehat\Omega_2) =: d$ für alle $x \in \partial\widehat\Omega_1$. Dann gilt aber auch $B_{\frac{d}{2}}(x) \subseteq \widehat\Omega_2$, also insbesondere $x \in \widehat\Omega_2$ für alle $x \in \partial\widehat\Omega_1$. Damit folgt 
  $$
  \widehat\Omega_1 \cup \partial\widehat\Omega_1 = \overline{\widehat\Omega}_1 \subseteq \widehat\Omega_2.
  $$

  Andererseits gilt $\varepsilon_2 < \dist(\Omega_1, \partial\Omega_2)$, denn
  \begin{align*}
  \varepsilon_2 
  &< \frac{1}{2} \dist(\Omega_1, \partial\widetilde\Omega_2) \\
  & \leq \frac{1}{2} (\dist(\Omega_1, \partial\widehat\Omega_2) + \dist(\partial\widehat\Omega_2, \widetilde\Omega_2)) \\
  & \leq \frac{1}{2} (\dist(\Omega_1, \partial\widehat\Omega_2) + \varepsilon_2).
  \end{align*}

  Setzt man das beschriebene Vorgehen induktiv fort, so erhält man eine Folge $(\Omega_j)_{j \in \N}$ von \lipschitz\hyp{}Gebieten und eine Folge $(\varepsilon_j)_{j \in \N}$ für die nach Konstruktion $0 < \varepsilon_j < \tfrac{1}{j}$ für alle $j \in \N$ gilt.
  Die Eigenschaften a), b) und c) werden also erfüllt.

  Es gilt noch zu zeigen, dass die so konstruierte Folge $(\Omega_j)_{j \in \N}$ auch Eigenschaft d) erfüllt, also $ \Omega \subseteq \bigcup_{j \in \N} \Omega_j $ gilt.
  Sei dazu $x \in \Omega$ beliebig.
  Weil $\Omega$ zusammenhängend ist, existiert ein $j_0 \in \N$, sodass
  $$
  x \in \widetilde\Omega_{j_0} \subseteq \Omega \cap B_{j_0}(x_0)
  $$
  gilt.
  Sei $d:=\dist(x, \partial\widetilde\Omega_{j_0})$.
  Dann existiert ein $j_1 > j_0$ mit $\varepsilon_{j_1} < d$.
  Da die Inklusion $\widetilde\Omega_{j_0} \subseteq \widetilde\Omega_{j_1}$ gilt, folgt mit Lemma \ref{lem:distBoundary} die Ungleichung $\varepsilon_{j_1} < \dist(x,\partial\widetilde\Omega_{j_1})$, was wiederum $x \in \widehat\Omega_{j_1} = \Omega_{j_1}$ impliziert.
  Damit gilt auch Eigenschaft d).
\end{proof}

\begin{bemnumber}
  \label{bem:boundedSubset}
  Mit der im Beweis von Lemma \ref{lem:lipschitzExhaust} verwendeten Konstruktion gilt nun auch, dass für alle beschränkten Teilgebiete $\Omega' \subseteq \Omega$ mit $\overline{\Omega'} \subseteq \Omega$ ein $j \in \N$ existiert, sodass $\Omega' \subseteq \Omega_j$
\end{bemnumber}

\section{Kompakte Einbettungen}

\begin{itemize}
  \item \cite{sohr2001navier}[S.58, Lemma 1.5.4]
\end{itemize}

\begin{lem}
  \label{lem:compactEmbedding0}
  Sei $\Omega \subseteq \R^n$ mit $n \geq 1$ ein beschränktes Gebiet und $1 < q < \infty$.
  Dann ist die Einbettung
  $$
  L^q(\Omega) \subseteq W^{-1,q}(\Omega)
  $$
  kompakt.
\end{lem}

\begin{lem}
  \label{lem:compactEmbedding}
  Sei $\Omega \subseteq \R^n$ mit $n \geq 2$ ein beschränktes \lipschitz\hyp{}Gebiet und $\Omega_0 \subseteq \Omega$ ein nichtleeres Teilgebiet.
  Zudem sei $1 < q < \infty$.
  Dann gilt die Ungleichung
  \begin{equation}
    \label{eq:compactInequality}
    \norm{u}_{L^q(\Omega)} 
    \leq C_1 \norm{\nabla u}_{W^{-1,q}(\Omega)^n}
    \leq C_1 C_2 \norm{u}_{L^q(\Omega)}
  \end{equation}
  für alle $u \in L^q(\Omega)$, welche die Integralgleichung
  $$
  \int_{\Omega_0} u \d x = 0
  $$
  erfüllen.
  Hierbei bezeichnen $C_1 = C_1(q, \Omega, \Omega_0) > 0$ sowie $C_2 = C_2(n) > 0$ Konstanten.
\end{lem}

\begin{proof}
  Wir halten zunächst fest, dass für alle Folgen $(v_j)_{j \in \N}$ mit $v_j \in C_0^\infty(\Omega)^n$ für alle $j \in \N$ und $\lim_{j \to \infty} \norm{v_j - v}_{1,q} = 0$ für ein $v \in \overline{C_0^\infty(\Omega)^n}^{\norm{\cdot}_{1,q}}$ auch $\lim_{j \to \infty} \norm{\div v_j - \div v} = 0$ gilt.  
  Hiermit ergibt sich für ein $u \in L^q(\Omega)$ folgende Darstellung der linearen Fortsetzung der Distribution $\nabla u \in (C_0^\infty(\Omega)^n)'$ auf den Raum $\overline{C_0^\infty(\Omega)^n}^{\norm{\cdot}_{1,q}}$:
  $$
    [\nabla u, \cdot] \colon v \to [\nabla u, v] 
    = - \langle u, \div v \rangle 
    = \int_{\Omega} u\; \div v \d x.
  $$
  Wir können somit $\nabla u$ als Element von $W^{-1,q}(\Omega)$ auffassen.
  Mit dieser Eigenschaft lässt sich der zweite Teil der Ungleichung (\ref{eq:compactInequality}) beweisen. 
  Es gilt nämlich für alle $v \in W_0^{1,q'}(\Omega)^n$
  \begin{align*}
    |[\nabla u, v]|
    &= | \langle u, \div v \rangle| \\
    &\leq \norm{u}_q \norm{\div v}_{q'} \tag{\hoelder}\\
    &\leq \sum_{|\alpha| = 1} \norm{u}_q \norm{D^\alpha v}_{q'} \\
    &\leq C_2 \norm{u}_q \norm{v}_{W^{1,q'}(\Omega)^n},
  \end{align*}
  mit einem aus der Normäquivalenz auf $\R^n$ stammenden $C_2 = C_2(n)$.

  Wir beweisen nun den ersten Teil der Ungleichung (\ref{eq:compactInequality}) durch einen Widerspruchsbeweis.
  Dazu nehmen wir an, es existiere keine Konstante $C > 0$, sodass die Ungleichung
  $$
  \norm{u}_q \leq C \norm{\nabla u}_{-1,q}
  $$
  für alle $u \in L^q(\Omega)$ mit $\int_{\Omega0} u \d x = 0$ gelte.
  Dann existiert insbesondere für alle $j \in \N$ ein $u_j \in L^q(\Omega)$ mit
  $$
  \norm{u_j}_q > j \norm{\nabla u_j}_{-1,q}
  $$
  und $\int_{\Omega_0} u_j \d x = 0$.
  Wir betrachten nun die normierte Folge $(\tilde u_j)_{n \in \N}$ mit 
  $$
  \tilde u_j := \frac{u_j}{\norm{u_j}_{q}}, \quad n \in \N.
  $$
  So gilt weiterhin aufgrund der Homogenität des Integrals $\int_{\Omega_0} \tilde u_j \d x = 0$. 
  Darüberhinaus gilt $\norm{\tilde u_j}_q = 1$ und die Ungleichung 
  \begin{displaymath}
    \norm{\nabla \tilde u_j}_{-1,q} < \frac{1}{j}, \quad j \in \N. \tag{$\ast$}
  \end{displaymath}

  Nach dem Satz von \banach\hyp{}\alaoglu\ ist die Einheitskugel in reflexiven \banach\hyp{}Räumen schwach kompakt.
  Damit folgt, dass die beschränkte Folge $(\tilde u_j)_{j \in \N}$ eine schwach konvergente Teilfolge mit Grenzwert $u \in L^q(\Omega)$ enthält.
  Zur Vereinfachung bezeichnen wir diese konvergente Teilfolge wieder mit $(\tilde u_j)_{j \in \N}$.
  Nach Definition der schwachen Konvergenz gilt somit
  $$
  \langle u, v \rangle = \lim_{j \to \infty} \langle \tilde u_j, v \rangle
  $$
  für alle $v \in L^{q'}(\Omega)$.
  Insbesondere gilt aufgrund der vorausgesetzten Beschränktheit von $\Omega$ auch $\1 \in L^q(\Omega)$ und damit
  $$
  \langle u, \1 \rangle = \int_{\Omega_0} u \d x = 0.
  $$
  Unter Verwendung der Ungleichung ($\ast$) folgt
  \begin{displaymath}
    \lim_{j \to \infty} \norm{\nabla \tilde u_j}_{-1,q} = 0 \tag{$\ast\ast$}.
  \end{displaymath}
  Für alle $v \in W_0^{1,q'}(\Omega)^n$ und $j \in \N$ gilt zudem 
  \begin{align*}
    |\nabla \tilde u_j, v]|
    &= |\langle \tilde u_j, \div v \rangle| \\
    &\leq \norm{\tilde u_j}_q \norm{\div v}_{1,q} \tag{\hoelder}\\
    &\leq \norm{\tilde u_j}_q \norm{v}_{1,q} \\
    &\leq \norm{\nabla \tilde u_j}_{-1,q} \norm{v}_{1,q}, \tag{\poincare}
  \end{align*}
  woraus schließlich
  \begin{align*}
    |[\nabla u, v]|
    &= |\langle u, \div v \rangle \\
    &= \lim_{j \to \infty}|\langle \tilde u_j, \div v \rangle \\
    &= \lim_{j \to \infty} |[\nabla \tilde u_j, v ]| \\
    &= 0
  \end{align*}
  folgt.
  Im distributionellen Sinne gilt damit $\nabla u = 0$.
  In Abschnitt \ref{subsec:smoothMolly} haben wir gezeigt, das dies gerade impliziert, dass $u$ konstant ist.
  Aus der Bedingung $\int_{\Omega_0} u \d x = 0$ folgt nun $u = 0$. 
  
  Nach \cite{sohr2001navier}[S.45 Lemma 1.1.3] existiert eine Konstante $C > 0$, sodass für alle $j \in \N$ die Ungleichung
  \begin{displaymath}
    1 = \norm{\tilde u_j}_q \leq C (\norm{\nabla \tilde u_j}_{-1,q} + \norm{\tilde u_j}_{-1,q}) \tag{$\ast\ast\ast$}
  \end{displaymath}
  gilt.
  Da die Folge $(\tilde u_j)_{j \in \N}$ beschränkt in $L^q(\Omega)$ ist und zudem die Einbettung nach Lemma $\ref{lem:compactEmbedding0}$ kompakt ist, existiert eine bezüglich der Norm auf $W^{-1,q}(\Omega)$ konvergente Teilfolge von $(\tilde u_j)_{j \in \N}$, welche gegen eine Funktion $\tilde u$ konvergiert. 
  Wir wollen die Teilfolge wieder mit $(\tilde u_j)_{j \in \N}$ bezeichnen.
  Insbesondere konvergiert die Folge $(\tilde u_j)_{j \in \N}$ schwach gegen $\tilde u$.
  Dann gilt aber $\lim_{j \to \infty} \tilde u_j = \tilde u = u = 0$ aufgrund der \hausdorff\hyp{}Eigenschaft der schwachen Topologie.

  Aus den Gleichungen ($\ast\ast$) und ($\ast\ast\ast$) folgt nun der Widerspruch
  \begin{displaymath}
  1 \leq \lim_{j \to \infty} ( \norm{\nabla \tilde u_j }_{-1,q} + \norm{\tilde u_j}_{-1,q} ) = 0.\qedhere
  \end{displaymath}
  

\end{proof}

\section{Darstellung von Funktionalen}

Zunächst beschäftigen wir uns mit dem Gradientenoperator.
Wir wollen zeigen, dass er unter gewissen Zusatzvoraussetzungen ein abgeschlossenes Bild besitzt.
Dies ist Inhalt des folgenden Lemmas.

\begin{lem}
  \label{lem:closedImageGradient}
  Sei $\Omega \subseteq \mathbb{R}^n$ mit $n \geq 2$ ein beschränktes Gebiet und $1 < q < \infty$.
  Dann gilt für die Abbildung
  $$
  \nabla \colon L^{q}(\Omega)^{n} \to (C_0^\infty(\Omega)^n)',
  $$
  dass das Bild
  $$
  \{\nabla v \in L^{q}(\Omega)^{n^2} \colon v \in W_0^{1,q}(\Omega)^n\} \subseteq L^{q}(\Omega)^{n^2}
  $$
  der Einschränkung von $\nabla$ auf $W_0^{1,q}$ eine abgeschlossene Teilmenge des Raumes $L^q(\Omega)^{n^2}$ ist.
\end{lem}

\begin{proof}
  Sei $(v_j)_{j \in \N}$ eine Folge in $W_0^{1,q}(\Omega)^n$, sodass die Folge der Gradienten $(\nabla v_j)_{j \in \N}$ in $L^q(\Omega)^{n^2}$ konvergiert.
  Als konvergente Folge ist diese insbesondere eine \cauchy\hyp{}Folge.
  Da $\Omega$ nach Voraussetzung ein beschränktes Gebiet ist, lässt sich die \poincare\hyp{}Ungleichung anwenden.
  Daraus folgt, dass auch die Folge $(v_j)_{j \in \N}$ eine \cauchy\hyp{}Folge in $L^q(\Omega)^n$ ist.
  Dann gilt jedoch mit der Definition der \sobolev\hyp{}Norm die Ungleichung
  $$
  \norm{v_j - v_k}_{W_{1,q}} \leq c (\norm{v_j - v_k}_q + \norm{\nabla v_j - \nabla v_k}_q)
  $$
  für ein aufgrund der verwendeten Normäquivalenz existierendes $c > 0$ und alle $j,k \in \N$.
  Die Folge $(v_j)_{j \in \N}$ ist also auch eine \cauchy\hyp{}Folge bezüglich \sobolev\hyp{}Norm.
  Der Raum $W_0^{1,q}$ ist nach Definition ein bezüglich \sobolev\hyp{}Norm abgeschlossener Unterraum von $W^{1,q}$ und damit gilt $\lim_{j \to \infty} v_j = v \in W_0^{1,q}$.
  Da zudem
  $$
  \norm{\nabla v_j}_{L_{1,q}} \leq \norm{v}_{W_{1,q}}  
  $$
  gilt, muss also auch $\lim_{j \to \infty} \nabla v_j = \nabla v$ gelten.
  Dies beweist, dass $D$ ein abgeschlossener Unterraum von $L^q(\Omega)^{n^2}$ ist.
\end{proof}

Basierend auf \cite{sohr2001navier}[S.61, Lemma 1.6.1] beweisen wir nun eine Verallgemeinerung dieses Lemmas. 

\begin{lem}
  \label{lem:divRepresentation}
  Sei $\Omega \subseteq \R^n$ mit $n \geq 2$ ein beschränktes Gebiet und $f \in W^{-1,q}(\Omega)^n$ mit $1 < q < \infty$.
  Dann existiert eine Matrix $F \in L^q(\Omega)^{n^2}$, welche die Gleichung
  $$ f = \div F $$
  im distributionellen Sinne und die Ungleichungen
  $$
  \norm{ f }_{W^{-1,q}(\Omega)^n} 
  \leq \norm{ F }_{L^q(\Omega)^{n^2}} 
  \leq C \norm{ f }_{W^{-1,q}(\Omega)^n}
  $$
  mit $C = C(\Omega) > 0$ erfüllt.
\end{lem}

\begin{proof}
  Wir betrachten den Raum
  $$
  D := \{\nabla v \in L^{q'}(\Omega)^{n^2} \colon v \in W_0^{1,q'}(\Omega)^n\} \subseteq L^{q'}(\Omega)^{n^2}
  $$
  der Gradienten $\nabla v = (D_j v_l)_{j,l=1}^n$ von Funktionen $v = (v_1,\dots,v_n) \in W_0^{1,q'}(\Omega)^n$.
  Nach Lemma \ref{lem:closedImageGradient} ist $D$ ein abgeschlossener Unterraum von $L^{q'}(\Omega)^{n^2}$.

  Wir definieren das Funktional
  $$
  \tilde f \colon \nabla v \mapsto [\tilde f, \nabla v]\;, \quad \nabla v \in D
  $$
  durch $[\tilde f, \nabla v] := [f, v]$ für alle $v \in W_0^{1,q'}(\Omega)^n$.
  Dann liefert die \hoelder\hyp{}Ungleichung zusammen mit der \poincare\hyp{}Ungleichung eine Konstante $C = C(\Omega) > 0$, sodass 
  $$
  |[\tilde f, \nabla v]| 
  = |[f, v]| 
  \underset{\hoelder}{\leq} \norm{f}_{-1,q} \norm{v}_{1,q'}
  \underset{\poincare}{\leq} C \norm{f}_{-1,q} \norm{\nabla v }_{q'}
  $$
  für alle $\nabla v \in D$ gilt.
  Somit ist $\tilde f$ ein stetiges Funktional auf $D \subseteq L^{q'}(\Omega)^{n^2}$ mit 
  $$\norm{\tilde f}_{D'} \leq C\norm{f}_{-1,q}.$$
  Der Satz von \hahn\hyp{}\banach\ liefert eine normgleiche Fortsetzung von $D$ nach $L^{q'}(\Omega)^{n^2}$.
  Nach dem Darstellungssatz von \riesz\ über Funktionale auf $L^{q}$  existiert nun eine Matrix $F \in L^{q}(\Omega)^{n^2}$ mit
  $$
  \langle F, \nabla v \rangle
  = \sum_{j,l=1}^n \int_\Omega F_{jl}(D_j v_l) \d x
  = \int_\Omega F \cdot \nabla v \d x
  = [\tilde f, \nabla v] 
  = [f, v]
  $$
  für alle $v = (v_1,\dots,v_n) \in W_0^{1,q'}(\Omega)^n$.
  Beweise des Darstellungssatzes finden sich in \cite{adams2003sobolev}[S.47, Theorem 2.44] und \cite{werner2011fa}[?].
  Zudem gilt
  $$ \norm{F}_{L^q(\Omega)^{n^2}} = \norm{\tilde f}_{(L^{q'}(\Omega)^{n^2})'} \leq C \norm{f}_{-1,q},$$
  da die Identifikation von Funktionalen auf $L^{q'}(\Omega)^{n^2}$ mit Funktionen aus $L^q(\Omega)^{n^2}$ isometrisch ist.
  Des Weiteren gilt für alle $v \in W_0^{1,q'}(\Omega)^n$ mit der \hoelder\hyp{}Ungleichung
  \begin{align*}
    |[f,v]| 
    &= |\langle F, \nabla v\rangle|  \\
    &\leq \norm{F}_q \norm{\nabla v}_{q'}  \tag{\hoelder\hyp{}Ungleichung}\\
    &\leq \norm{F}_q (\norm{v}_{q'}^{q'} + \norm{\nabla v}_{q'}^{q'})^{\frac{1}{q'}} \tag{Übergang zu \sobolev\hyp{}Norm}\\
    &= \norm{F}_q \norm{v}_{W_0^{1,q'}(\Omega)^n}.
  \end{align*}
  Daraus folgt
  $$
  \norm{f}_{W^{-1,2}(\Omega)} \leq \norm{F}_2.
  $$
  Ist nun $v \in C_0^\infty(\Omega)^n$, so gilt
  \begin{align*}
    [f,v]
    &= \langle F, \nabla v\rangle \\
    &= \sum_{j,l=1}^n \langle F_{jl}, D_j v_l \rangle \\
    &= -\sum_{j,l=1}^n \langle D_j F_{jl}, v_l \rangle \\
    &= -[\div F, v], \\
  \end{align*}
  die Abbildungen $[f,\cdot]$ und $[-\div F, \cdot]$ stimmen also im Sinne von Distributionen überein, die zu zeigende Aussage folgt durch den Übergang von $F$ zu $-F$.
\end{proof}

\section{Das Gradientenkriterium}

In diesem Abschnitt stellen wir ein für die \helmholtz\hyp{}Zerlegung fundamentales Kriterium vor, welches es ermöglicht, Funktionen $f \in W^{1,q}(\Omega)^n$ als Gradienten $f = \nabla p$ mit $p \in L^q(\Omega)$ darzustellen.

\begin{lem}
  Sei $\Omega \subseteq \R^n$ mit $n \geq 2$ ein Gebiet $\Omega_0 \subseteq \Omega$ ein beschränktes Teilgebiet mit $\emptyset \neq \overline\Omega_0 \subseteq \Omega$ und $1 < q < \infty$.
  Angenommen, für $f \in W_{\loc}^{-1,q}(\Omega)^n$ gelte
  \begin{equation}
    \label{eq:vKernel}
    [f,v] = 0, \quad \text{für alle} \quad v \in C_{0,\sigma}^\infty(\Omega).
  \end{equation}
  Dann existiert ein eindeutig bestimmtes $p \in L^q_{\loc}(\Omega)$, welches die Gleichung $\nabla p = f$ im distributionellen Sinne erfüllt und für das zusätzlich
  \begin{equation}
    \int_{\Omega_0} p \d x = 0
    \label{eq:0integral}
  \end{equation}
  gilt.
\end{lem}

\begin{proof}
  Wir zeigen zunächst, dass für ein beliebiges beschränktes \lipschitz\hyp{}Gebiet $\Omega_1 \subseteq \Omega$ mit $\overline\Omega_0 \subseteq \Omega_1 \subseteq \overline\Omega_1 \subseteq \Omega$ ein eindeutig bestimmtes $p \in L^q(\Omega_1)$ existiert, welches die Behauptung des Lemmas erfüllt.

  Ähnlich zum ersten Beweisschritt von Lemma \ref{lem:lipschitzExhaust} finden wir ein weiteres beschränktes \lipschitz\hyp{}Gebiet $\Omega_2$ mit $\overline\Omega_1 \subseteq \Omega_2 \subseteq \overline\Omega_2 \subseteq \Omega$.
  Dazu wählen wir ein $x_0 \in \Omega_1$ und finden aufgrund der vorausgesetzten Beschränktheit von $\Omega_1$ ein $r > 0$, sodass $\Omega_1 \subseteq B_r(x_0)$ gilt.
  Wir wählen sodann die $x_0$ enthaltende Zusammenhangskomponente $\widetilde\Omega_2$ von $B_r(x_0) \cap \Omega$ aus und konstruieren wie schon im Beweis von Lemma \ref{lem:lipschitzExhaust} das beschränkte \lipschitz\hyp{}Gebiet $\Omega_2 = \widehat\Omega_2$.

  Der Voraussetzung $f \in W_{\loc}^{-1,q}(\Omega)^n$ entnehmen wir, dass $f \in W^{-1,q}(\Omega_2)^n$ gilt. Zudem existiert aufgrund der Beschränktheit von $\Omega_2$ nach Lemma \ref{lem:divRepresentation} eine Darstellung
  $$
  f = \div F \quad\text{mit}\quad F = (F_{jl})_{j,l=1}^n \in L^q(\Omega_2)^{n^2}.
  $$

  Im Folgenden bezeichne $F^\varepsilon := \mathcal{F}_\varepsilon \ast F = (\mathcal{F}_\varepsilon \ast F_{jl})_{j,l=1}^n$ mit $0 < \varepsilon < \dist(\Omega_1, \partial\Omega_2)$ die in Abschnitt \ref{subsec:smoothMolly} definierte Faltung von $F$ mit einem Glättungskern, für den wie bereits gezeigt $F^\varepsilon \in C^\infty(\overline\Omega_1)^{n^2}$ gilt.
  Wir wollen beweisen, dass eine Darstellung der Form
  $$
  \div F^\varepsilon = \nabla U_\varepsilon
  $$
  mit einem $U_\varepsilon \in C^\infty(\overline\Omega_1)$ existiert.

  Ein bekanntes Resultat aus der Funktionentheorie besagt, dass jedes auf dem Gebiet $\Omega_1$ definierte glatte Vektorfeld $g \in C^\infty(\Omega_1)^n$ genau dann ein Potential $\Phi \in C^\infty(\Omega_1)$ besitzt, falls das Kurvenintegral entlang jeder in $\Omega_1$ verlaufenden glatten Kurve $w \colon [0,1] \to \Omega_1$ verschwindet, also genau dann, falls
  $$
  \oint_w g \cdot \d s = \int_0^1 g(w(\tau))\cdot w'(\tau) \d\tau = 0
  $$
  für alle glatten Kurven $w$ in $\Omega_1$ gilt.
  %Da sich stetige Funktionen und damit auch die Ableitungen glatter Funktionen eindeutig und stetig auf den Rand ihres Definitionsbereichs fortsetzen lassen, kann man auch $g \in C^\infty(\overline\Omega_1)^n$ sowie $\Phi \in C^\infty(\overline\Omega_1)$ annehmen.

  Wir wollen dies nun für $\div F^\varepsilon$ beweisen.
  Dazu definieren wir für alle $x\in\Omega_2$ den Wert des Integrals
  $$
  V_{w,\varepsilon}(x) := \int_0^1 \mathcal{F}_\varepsilon(x - w(\tau))w'(\tau)\d\tau
  $$
  und erhalten durch wiederholte Anwendung der \leibniz\hyp{}Regel für Parameterintegrale $ V_{w,\varepsilon} \in C_0^\infty(\Omega_2)^n$.
  Darüber hinaus gilt für eine geschlossene Kurve $w$ in $\overline\Omega_1$ folgende Rechnung
  \begingroup
  \addtolength{\jot}{0.5em}
  \begin{align*}
    \div V_{w,\varepsilon}(x) 
    &= \int_0^1 \sum_{j=1}^n (D_j \mathcal{F}_\varepsilon)(x - w(\tau))w_j'(\tau) \d\tau \tag{\leibniz}\\
    &= -\int_0^1 \frac{\d\ }{\d\tau}\mathcal{F}_\varepsilon(x - w(\tau)) \d\tau \tag{Kettenregel}\\
    &= \mathcal{F}_\varepsilon(x - w(0)) - \mathcal{F}_\varepsilon(x - w(1)) \tag{Hauptsatz}\\
    &= 0 \tag{geschlossene Kurve}
  \end{align*}
  \endgroup

  Hiermit folgt $V_{w,\varepsilon} \in C_{0,\sigma}^\infty(\Omega_2)^n$.
  Unter Verwendung der Voraussetzung aus Gleichung (\ref{eq:vKernel}) und dem Satz von \fubini\ folgt
  \begingroup
  \addtolength{\jot}{1em}
  \begin{align*}
    0
    %
    &= [f, V_{w, \varepsilon}] %\qquad &&\text{Voraussetzung (\ref{eq:vKernel})}\\
    %
    = [\div F, V_{w, \varepsilon}] %\\
    %
    = \int_{\Omega_2} \div F \cdot V_{w, \varepsilon} \d x \\
    %
    &= \int_{\Omega_2} \sum_{j,l=1}^n D_j F_{jl}(x) \left( \int_0^1 \mathcal{F}_\varepsilon(x - w(\tau))w_l'(\tau) \d\tau \right) \d x \\
    %
    %&= \sum_{j,l=1}^n \int_{\Omega_2} D_j F_{jl}(x) \left( \int_0^1 \mathcal{F}_\varepsilon(x - w(\tau))w_l'(\tau) \d\tau \right) \d x \\
    %
    %&= \sum_{j,l=1}^n \int_{\Omega_2} \int_0^1 D_j F_{jl}(x) \mathcal{F}_\varepsilon(x - w(\tau))w_l'(\tau) \d\tau \d x \\
    %
    %&= \sum_{j,l=1}^n \int_0^1 \int_{\Omega_2}  D_j F_{jl}(x) \mathcal{F}_\varepsilon(x - w(\tau))w_l'(\tau) \d\tau \d x 
    %
    &= \sum_{j,l=1}^n \int_0^1 \left(\int_{\Omega_2}  D_j F_{jl}(x) \mathcal{F}_\varepsilon(x - w(\tau))\d x \right) w_l'(\tau) \d\tau \tag{Satz von \fubini}\\
    %
    &= \sum_{j,l=1}^n \int_0^1 \left(\int_{\Omega_2}  D_j F_{jl}(x) \mathcal{F}_\varepsilon(w(\tau) - x)\d x \right) w_l'(\tau) \d\tau \tag{$\mathcal{F}_\varepsilon(x) = \mathcal{F}_\varepsilon(-x)$}\\
    %
    &= \sum_{j,l=1}^n \int_0^1 \left(D_j F_{jl} \ast \mathcal{F}_\varepsilon\right)(w(\tau)) w_l'(\tau) \d\tau \\
    %
    &= \sum_{j,l=1}^n \int_0^1 D_j\left(F_{jl} \ast \mathcal{F}_\varepsilon\right)(w(\tau)) w_l'(\tau) \d\tau \tag{1.7.17}\\
    %
    &= \sum_{j,l=1}^n \int_0^1 (D_j F_{jl}^\varepsilon)(w(\tau)) w_l'(\tau) \d\tau \\
    %
    &= \int_0^1 (\div F^\varepsilon)(w(\tau)) \cdot w'(\tau) \d\tau \\
    %
    &= \oint_w \div F^\varepsilon \cdot \d s.
  \end{align*}
  \endgroup

  Es existiert also ein $U_\varepsilon \in C^\infty(\overline\Omega_1)$ mit $\div F^\varepsilon = \nabla U_\varepsilon$.
  Die Funktion $U_\varepsilon$ ist bis auf eine Konstante $\gamma$ eindeutig bestimmt, denn es gilt $\nabla U_\varepsilon = \nabla( U_\varepsilon - \gamma)$.
  Wir können durch die Wahl $\gamma = \int_{\Omega_0} U_\varepsilon \d x$ erreichen, dass $\int_{\Omega_0} U_\varepsilon - \gamma \d x = 0$  gilt.
  Im Folgenden bezeichnen wir die so gewählte Funktion wieder mit $U_\varepsilon$.
  Mit Lemma \ref{lem:compactEmbedding} folgt
  \begingroup
  \addtolength{\jot}{1em}
  \begin{align*}
    \norm{U_\varepsilon}_{L^q(\Omega_1)}
    &\leq C \norm{ \nabla U_\varepsilon }_{W^{-1,q}(\Omega_1)} \\
    &= C \sup_{0 \neq v \in C_0^\infty(\Omega_1)^n} \left(\frac{|[\nabla U_\varepsilon, v]|}{\norm{\nabla v}_{q'}} \right) \\
    &\overset{(\ast)}{=} C \sup_{0 \neq v \in C_0^\infty(\Omega_1)^n} \left( \frac{|\langle F^\varepsilon, \nabla v \rangle|}{\norm{v}_{q'}}\right) \\
    &\leq C \norm{F^\varepsilon}_{L^q(\Omega_1} \tag{$\star$}
  \end{align*}
  \endgroup
  für ein $C = C(q,\Omega_0,\Omega_1)$, welches nicht von $\varepsilon$ abhängt. 
  Die Gleichheit ($\ast$) ergibt sich dabei wie folgt:
  \begin{align*}
    \langle F^\varepsilon, \nabla v \rangle
    &= \int_{\Omega_1} F^\varepsilon \cdot \nabla v \d x \\
    &= \int_{\Omega_1} \sum_{j,l=1}^n F_{jl}^\varepsilon (D_j v_l) \d x \\
    &= \int_{\Omega_1} \sum_{j,l=1}^n -(D_jF_{jl}^\varepsilon) v_l \d x \tag{Kompakter Träger}\\
    &= \int_{\Omega_1} -\div F^\varepsilon \cdot v \d x \\
    &= \int_{\Omega_1} -\nabla U_\varepsilon \cdot v \d x \\
    &= -[\nabla U_\varepsilon, v].
  \end{align*}

  Eine wesentliche Eigenschaft der Glättungskerne ist, dass $\lim_{\varepsilon \to 0} \norm{F - F^\varepsilon}_{L^q(\Omega_1)} = 0$ gilt.
  Das Netz $(F^\varepsilon)_{\varepsilon \in \R^+}$ ist aufgrund der Konvergenz auch ein \cauchy-Netz.
  Mit Gleichung ($\star$) gilt dann für alle $0 < \eta < \varepsilon$ die Ungleichung
  $$
    \norm{U_\varepsilon - U_\eta}_{L^q(\Omega_1)}
    \leq C \norm{F^\varepsilon - F^\eta}_{L^q(\Omega_1}.
  $$
  Also ist auch $(U_\varepsilon)_{\varepsilon \in \R^+}$ ein \cauchy-Netz, welches aufgrund der Vollständigkeit des Raumes $L^q(\Omega_1)$ einen eindeutig bestimmten Grenzwert $U \in L^q(\Omega_1)$ besitzt.
  Da zudem $\Omega_0$ beschränkt ist, gilt die Inklusionsbeziehung $L^1(\Omega_0) \subseteq L^q(\Omega_0)$.
  Damit folgt
  $$
    | \int_{\Omega_0} U - U_\varepsilon \d x |
    \leq  \int_{\Omega_0} | U - U_\varepsilon | \d x 
    = \norm{U - U_\varepsilon}_{L^1(\Omega_0)} 
    \to 0 \quad\text{für } \varepsilon \to 0,
  $$
  also auch
  $$
    \int_{\Omega_0} U \d x = \lim_{\varepsilon \to 0} \int_{\Omega_0} U_\varepsilon \d x =  0,
  $$
  aufgrund der Wahl von $U^\varepsilon$.

  Nun zeigen wir noch dass $\div F = \nabla U$ im distributionellen Sinne gilt. 
  Wir halten dazu zunächst fest, dass die Konvergenz auf $L^q(\Omega_1)^n$ die komponentenweise Konvergenz auf $L^q(\Omega_1)$ impliziert.
  Es gilt somit $\norm{F_{jl} - F_{jl}^\varepsilon}_{L^q(\Omega_1)} \to 0$ für $\varepsilon \to 0$.
  Sei nun $v \in C_0^\infty(\Omega_1)^n$.
  Dann gilt
  \begin{align*}
    [\nabla U, v]
    &= \int_{\Omega_1} \nabla U \cdot v \d x \\
    &= \int_{\Omega_1} \sum_{j = 1}^n (D_j U) v_j \d x \\
    &= -\int_{\Omega_1} \sum_{j = 1}^n U (D_j v_j) \d x \\ 
    &= -\int_{\Omega_1} \sum_{j = 1}^n \lim_{\varepsilon \to 0} U_\varepsilon (D_j v_j) \d x \\
    &\overset{(\clubsuit)}{=} -\lim_{\varepsilon \to 0} \int_{\Omega_1} \nabla U_\varepsilon \cdot v \d x \\
    &= -\lim_{\varepsilon \to 0} \int_{\Omega_1} \div F^\varepsilon \cdot v \d x \\
    &= -\lim_{\varepsilon \to 0} \int_{\Omega_1} \sum_{j,l=1}^n F_{jl}^\varepsilon (D_j v_l) \d x \\
    &\overset{(\clubsuit)}{=} \int_{\Omega_1} \sum_{j,l=1}^n (D_j F_{jl}) v_l \d x \\
    &= \langle \div F, v \rangle,
  \end{align*}
  wobei bei ($\clubsuit$) verwendet wurde, dass $\Omega_1$ beschränkt ist.

  Zuletzt zeigen wir, wie sich die gesuchte Funktion $p$ konstruieren lässt.
  Wie der Beweis zeigt, lässt sich für jedes beschränkte \lipschitz\hyp{}Gebiet $\Omega_1$ ein $U \in L^q(\Omega_1)$ finden, welches zudem aufgrund der Forderung 
  $$
  \int_{\Omega_0} U \d x = 0
  $$
  eindeutig bestimmt ist.
  Nach Lemma \ref{lem:lipschitzExhaust} lässt sich jedes Gebiet $\Omega$ durch eine Folge $(\Omega_j)_{j \in \N}$ beschränkter \lipschitz\hyp{}Gebiete ausschöpfen.
  Zudem gilt nach Bemerkung \ref{bem:boundedSubset} jedes beschränkte Teilgebiet $\Omega' \subseteq \Omega$ in einem beschränkten \lipschitz\hyp{}Gebiet $\Omega_j, j \in \N$ enthalten ist.
  Wir können direkt annehmen, dass dies bereits für $j = 1$ erfüllt ist.
  Führen wir nun den vorangehenden Beweis für $\Omega_j$, so erhalten wir eine eindeutig bestimmte Funktion $p \in L^q(\Omega_j)$, die die Gleichung $f = \nabla p$ auf $\Omega_j$ im distributionellen Sinne erfüllt und für die des Weiteren das Integral $\int_{\Omega_0} p \d x$ verschwindet.
  Da nach Lemma \ref{lem:lipschitzExhaust} die Gleichheit $\bigcup_{j \in \N} \Omega_j = \Omega$ gilt, lässt sich $p$ auch eindeutig bis auf Nullmengen auf $\Omega$ definieren.
  Ist zudem $B$ ein Ball mit $\overline B \subseteq \Omega$, so existiert ein $j \in \N$ mit $\overline B \subseteq \Omega_j$.
  Da $p \in L^q(\Omega_j)$ gilt, folgt sofort $p \in L^q(B)$.
  Die Definition der lokalen Integrierbarkeit liefert hiermit $p \in L^q_{\loc}(\Omega)$.
  Dies beweist die Behauptung.
\end{proof}

\chapter{Helmholtz-Zerlegung in $L^2$}

Nach der in den vorangehenden Kapiteln geleisteten Vorarbeit sind wir nun in der Lage den zentralen Satz dieser Arbeit zu behandeln.
Wir interessieren uns dafür, unter welchen Bedingungen sich die Vektorräume $L^q(\Omega)^n$ in für die Strömungsmechanik interessante Teilräume zerlegen lassen.
Die Darstellung in dieser Arbeit begrenzt sich auf die Zerlegung für $q = 2$.
Hier ermöglicht der für die Theorie der \hilbert\hyp{}Räume grundlegende Begriff der Orthogonalität eine orthogonale Zerlegung des Lösungsraumes $L^2(\Omega)^n$.

Wir bezeichnen mit 
$$
G(\Omega) := \{f \in L^2(\Omega)^n \mid f = \nabla p \text{ für ein } p \in L^2_{\loc}(\Omega)\}
$$
den Raum aller Funktionen $f \in L^2(\Omega)^n$, die im distributionellen Sinne ein skalares Potential besitzen.

\begin{thm}
  \label{lem:helmholtz}
  Sei $\Omega \subseteq \R^n$ mit $n \geq 2$ ein beliebiges Gebiet.
  Dann gilt
  $$
  L^2(\Omega)^n = G(\Omega) \oplus L^2_\sigma(\Omega).
  $$
\end{thm}

\begin{proof}
  Wir zeigen, dass 
  \begin{equation}
    \label{eq:orthComplement}
    G(\Omega) = L_\sigma^2(\Omega)^\perp
  \end{equation}
  gilt.
  Die zu zeigende Aussage folgt dann aus dem Fundament der Theorie der \hilbert\hyp{}Räume: dem Satz von der Orthogonalprojektion.

  Sei dazu $f \in L_\sigma^2(\Omega)^\perp$.
  Wir fassen $f$ als Distribution auf, also als Element von $(C_0^\infty(\Omega)^n)'$.
  Für jedes beschränkte Teilgebiet $\Omega_0 \subseteq \Omega$ mit $\overline\Omega_0 \subseteq \Omega$ folgt für alle $v \in C_0^\infty(\Omega_0)^n$ mit der \cauchy\hyp\schwartz-\ (C.S.) und \poincare\hyp{}Ungleichung \ref{lem:poincare}
  $$
  |[f, v]| 
  = |\langle f, v \rangle| 
  \underset{\textsc{C.S.}}{\leq}  \norm{f}_{L^2(\Omega_0)^n} \norm{v}_{L^2(\Omega_0)^n}
  \underset{\poincare}{\leq} C \norm{f}_{L^2(\Omega_0)^n} \norm{\nabla v}_{L^2(\Omega_0)^{n^2}},
  $$
  mit einer Konstante $C = C(\Omega_0) > 0$.
  Es gilt also 
  \begin{align*}
    |[f, v]| 
    &\leq \norm{f}_{L^2(\Omega_0)^n} \cdot ( C + 1) \cdot (\norm{v}_{L^2(\Omega_0)^n} + \norm{\nabla v}_{L^2(\Omega_0)^{n^2}}) \\
    &\leq \norm{f}_{L^2(\Omega_0)^n} \cdot C'(n) \cdot ( C + 1) \cdot \norm{v}_{1,2}
  \end{align*}
  mit einem aus der Normäquivalenz der Normen auf $\R^n$ folgenden $C'(n)$.
  Unter Berücksichtigung von Gleichung (\ref{eq:operatornorm}) folgt daraus $f \in W_{\loc}^{-1,2}(\Omega)^n$.
  Nach Voraussetzung gilt zudem 
  $$
  [f,v] = \langle f, v \rangle = 0
  $$
  für alle $v \in C_{0, \sigma}^\infty(\Omega)$.
  Die Anwendung des Gradientenkriteriums aus Lemma \ref{lem:gradientCriterion} liefert nun die Existenz einer Funktion $p \in L_{\loc}^2(\Omega)$, welche die Gleichung
  $$
  f = \nabla p
  $$
  im distributionellen Sinne löst.
  Daraus folgt $f \in G(\Omega)$.

  Sei umgekehrt $f \in G(\Omega)$ mit $f = \nabla p$ für ein $p \in L^2_{\loc}(\Omega)$.
  Dann gilt
  \begin{align*}
  \langle f, v \rangle
  &= [ \nabla p, v ] \\
  &= \sum_{j = 1}^n [ D_j p , v_j ] \\
  &= - \sum_{j = 1}^n \langle p, D_j v \rangle \\
  &= - \langle p, \div v \rangle  \\
  &= - \langle p, 0 \rangle  \\
  &= 0
  \end{align*}
  für alle $v \in C_{0, \sigma}^\infty(\Omega)$ und wegen der Beziehung $\nabla p \in L^2(\Omega)$ gilt dies aufgrund der Stetigkeit des Skalarproduktes auch für alle $v \in \overline{C_{0, \sigma}^\infty(\Omega)}_{\norm{\cdot}_2} = L_\sigma^2(\Omega)$.
  Damit ist $f \in L_{0,\sigma}^2(\Omega)^\perp$.
\end{proof}

Aus Gleichung (\ref{eq:orthComplement}) folgt insbesondere, dass der Raum $G(\Omega)$ bezüglich der Norm $\norm{\cdot}_2$ abgeschlossen ist.
Wie aus der Theorie der \hilbert\hyp{}Räume bekannt ist, existiert zu dem abgeschlossenen Teilraum $L_\sigma^2(\Omega)$ ein eindeutig bestimmter Projektionsoperator, die sogenannte \helmholtz\hyp{}Projektion
$$
P \colon L^2(\Omega)^n \to L_\sigma^2(\Omega),\quad Pf \mapsto f_0,
$$
für alle $f = f_0 + \nabla p \in L^2(\Omega)$ mit $f_0 \in L_\sigma^2(\Omega)$ und $\nabla p \in G(\Omega)$.

Im Falle $\Omega = \R^n$ lassen sich die Räume $G(\Omega)$ und $L_\sigma^2(\Omega)$ noch auf eine andere Art charakterisieren.
Dies ist Inhalt des folgenden Satzes.

\begin{thm}
  Sei $n \in \N$ mit $n \geq 2$.
  Dann gilt
  \begin{equation}
    \label{eq:LCharacterization}
    L_\sigma^2(\R^n) = \{ f \in L^2(\R^n)^n \mid \div f = 0 \}
  \end{equation}
  und $G(\R^n)$ ist der Abschluss des Raumes
  $$
  \nabla C_0^\infty(\R^n) := \{ \nabla p \mid p \in C_0^\infty(\R^n) \}
  $$
  bezüglich der Norm $\norm{\cdot}_{L^2(\R^n)^n}$.
  Es gilt also
  \begin{equation}
    \label{eq:GCharacterization}
    G(\R^n) = \overline{\nabla C_0^\infty(\R^n)}^{\norm{\cdot}_2}.
  \end{equation}
\end{thm}

\begin{proof}
  Im folgenden Beweis wollen wir eine Funktion $\varphi \in C_0^\infty(\R^n)$ mit den Eigenschaften
  \begin{equation}
    \label{eq:propPhi}
  0 \leq \varphi \leq 1
  \quad,\quad
  \varphi(x) = 1 \text{, falls } |x| \leq 1
  \quad,\quad
  \varphi(x) = 0 \text{, falls } |x| \geq 2,
\end{equation}
  betrachten.
  Eine solche Funktion erhält man zum Beispiel aus der Faltung der charakteristischen Funktion $\chi_{B_{3/2}}(0)$ mit dem Standardglättungskern $\mathcal{F}_{1/2}$.
  Zudem definieren wir die Funktionen
  $$
  \varphi_j \in C_0^\infty(\R^n)
  \quad,\quad
  \varphi_j(x) := \varphi\left(\frac{x}{j}\right)
  \quad,\quad
  x \in \R^n, j \in \N.
  $$
  Nach Konstruktion folgt 
  \begin{equation}
    \label{eq:limPhi}
    \lim_{j \to \infty} = \varphi_j(x) = 1
  \end{equation}
  für alle $x \in \R^n$.
  Setzen wir
  $$
  B_j := \{ x \in \R^n \mid |x| < j \} 
  \quad,\quad
  G_j := B_{2j}\setminus \overline B_j,
  $$
  so erhalten wir $\supp \nabla \varphi_j \subseteq \overline G_j$, da nach Definition $\varphi_j$ auf $B_j$ konstant ist sowie $\supp \varphi_j \subseteq \overline B_{2j}, j \in \N$ gilt.

  Wir beginnen damit, die Gültigkeit von Gleichung (\ref{eq:GCharacterization}) zu zeigen.
  Da sich der Raum $C_0^\infty(\R^n)$ als Unterraum von $L^2(\R^n)$ auffassen lässt, folgt sogleich $G(\R^n) \supseteq \nabla C_0^\infty(\R^n)$.
  Zudem folgt aus Satz \ref{lem:helmholtz}, dass $G(\R^n)$ als orthogonales Komplement ein bezüglich $\norm{\cdot}_2$ abgeschlossener Unterraum von $L^2(\R^n)^n$ ist.
  Damit gilt auch $G(\R^n) = \overline{G(\R^n)}^{\norm{\cdot}_2} \supseteq \overline{\nabla C_0^\infty(\R^n)}^{\norm{\cdot}_2}$.

  Für die andere Inklusion betrachten wir den Gradienten $\nabla p \in G(\R^n)$.
  Wir wählen Konstanten 
  $$
  K_j := \frac{\int_{G_j} p \d x}{\int_{G_j} \1 \d x}, \quad j \in \N.
  $$
  Aus der \poincare\hyp{}Ungleichung \cite[S.44, Lemma 1.1.2]{sohr2001navier} folgt sodann
  \begin{align}
  \norm{p - K_1}_{L^2(G_1)} &\leq C \left(\norm{\nabla (p - K_1)}_{L^2(G_1)^n} + | \int_{G_1} p - K_1 \d x |\right) \nonumber\\
    \label{eq:poincareP}
  &=  C \;\norm{\nabla p}_{L^2(G_1)^n}.
\end{align}
  Mit der Substitutionsregel erhalten wir
\begin{align*}
  \norm{p - K_j}_{L^2(G_j)} 
  &= \left( \int_{G_j} |p(x) - K_j|^2 \d x \right)^\frac{1}{2} \\
  &= \left( \int_{G_1} |p(jy) - K_j|^2 \d y \right)^\frac{1}{2} j^\frac{n}{2} \tag{Substitutionsregel}\\
  &\leq C j^{\frac{n}{2}}\left( \int_{G_1} |\nabla_y p(jy) |^2 \d y \right)^{\frac{1}{2}} \tag{mit Gleichung (\ref{eq:poincareP})}\\
  &= C j^{\frac{n}{2}}j^{-\frac{n}{2}}j \left( \int_{G_j} |\nabla p(x) |^2 \d x \right)^{\frac{1}{2}} \tag{Ketten- u. Substitutionsregel} \\
  &= C j \norm{\nabla p}_{L^2(G_j)^n}
\end{align*}
  für alle $j \in \N$.
  Zusammenfassend gilt also
  \begin{equation}
    \label{eq:gradientInequality}
    \norm{p - K_j}_{L^2(G_j)} \leq j C\norm{\nabla p}_{L^2(G_j)^n}.
  \end{equation}

  Wir definieren nun $p_j := \varphi_j (p - K_j)$.
  Die komponentenweise Anwendung der Produktregel ergibt 
  \begin{equation}
    \label{eq:productrule}
    \nabla p_j =  (\nabla \varphi_j)(p - K_j) + \varphi_j \nabla p,
  \end{equation}
  wobei $\nabla \varphi_j(x) = \frac{1}{j} \nabla \varphi (\frac{x}{j})$ gilt.
  Hieraus erhalten wir unter Verwendung der Konstanten $C' := \sup_x |\nabla \varphi(x)|$ und Gleichung (\ref{eq:gradientInequality}) die folgende Abschätzung:
  \begin{align}
    \norm{\nabla p - \nabla p_j}_{L^2(\R^n)^n}
    &\leq \norm{\nabla p - \varphi_j \nabla p}_{L^2(\R^n)^n} + \norm{\varphi_j \nabla p - \nabla p_j}_{L^2(\R^n)^n} \nonumber\\
    &\hspace{-0.15cm}\underset{(\ref{eq:productrule})}{=} \norm{\nabla p - \varphi_j \nabla p}_{L^2(\R^n)^n}    + \norm{\nabla\varphi_j (p - K_j)}_{L^2(\R^n)^n} \nonumber\\
    &\hspace{-0.15cm}\underset{(\ref{eq:gradientInequality})} \leq \norm{\nabla p - \varphi_j \nabla p}_{L^2(\R^n)^n} + \frac{C'}{j} \norm{p - K_j}_{L^2(G_j)^n}\nonumber\\
    &\leq \norm{\nabla p - \varphi_j \nabla p}_{L^2(\R^n)^n} + C'C \norm{\nabla p}_{L^2(G_j)^n} \label{eq:gradDiff}.
  \end{align}
  Hierbei wurde beim Übergang von $\norm{\cdot}_{L^2(\R^n)^n}$ zu $\norm{\cdot}_{L^2(G_j)^n}$ ausgenutzt, dass nach Konstruktion $\supp \nabla \varphi_j \subseteq G_j$ gilt.

  Aus Gleichung (\ref{eq:propPhi}) folgt $|1 - \varphi_j(x)| = |1 - \varphi(\tfrac{x}{j})| \leq 2$. 
  Die Glieder der Funktionenfolge $\left((1 - \varphi_j) D_k p\right)_{j \in \N}$ besitzen somit die integrierbare Majorante $2\, D_k p \in L^2(\R^n)$ für $k = 1, \dots n$.
  Mit dem Satz von \lebesgue\ über majorisierte Konvergenz ergibt sich daraus unter Verwendung von Gleichung (\ref{eq:limPhi})
  $$
  \lim_{j \to \infty} \norm{(1 - \varphi_j) \,D_k p}_2 
  = \left( \lim_{j \to \infty} |(1 - \varphi_j)|^2 |\,D_k p|^2 \d x \right)^{\frac{1}{2}}
  = 0,
  $$
  was wiederum
  \begin{equation}
    \label{eq:lebesgueDom}
  \lim_{j \to \infty} \norm{\nabla p - \varphi_j \nabla p}_{L^2(\R^n)^n}
  = \left(\, \sum_{k = 1}^n \lim_{j \to \infty} \norm{(1 - \varphi_j)\, D_k p}_2^2  \right)^\frac{1}{2} = 0
  \end{equation}
  impliziert.
  Ebenso folgt aus dem Satz über majorisierte Konvergenz 
  \begin{equation}
    \label{eq:limitGradient}
  \lim_{j \to \infty} \norm{\nabla p}_{L^2(G_j)^n}
  = 0.
  \end{equation}
  Aus Gleichung (\ref{eq:gradDiff}) folgt nun mit (\ref{eq:lebesgueDom}) und (\ref{eq:limitGradient})
  \begin{equation}
    \label{eq:limNablaP}
    \lim_{j \to \infty} \norm{\nabla p - \nabla p_j}_{L^2(\R^n)^n} = 0.
  \end{equation}

  Nach Konstruktion gilt $\nabla p_j \in L^2(\R^n)^n$.
  Durch Lemma \ref{lem:mollification} lässt sich somit $\nabla p_j$ als Grenzwert von glatten Funktionen mit kompaktem Träger darstellen. Für alle $j \in \N$ existiert daher ein $\varepsilon_j > 0$ mit
  \begin{equation}
    \label{eq:convolutionNablaP}
  \norm{\nabla p_j - \nabla (\mathcal{F}_{\varepsilon_j} \ast p_j)}_{L^2(\R^n)^n}
  \overset{(\ref{eq:convolutionDiff})}{=} \norm{\nabla p_j - \mathcal{F}_{\varepsilon_j} \ast \nabla p_j}_{L^2(\R^n)^n}
  \leq \frac{1}{j}.
  \end{equation}
  Definieren wir nun $\tilde p_j := \mathcal{F}_{\varepsilon_j} \ast p_j$ so gilt $\tilde p_j \in C_0^\infty$, da 
  $$
  \supp p_j = \supp \varphi_j(p - K_j) 
  \subseteq \supp \varphi_j
  \subseteq \overline B_{2j}
  $$
  impliziert, dass $p_j$ einen kompakten Träger besitzt und selbiges auch für $\tilde p_j$ als Faltung zweier Funktionen mit kompaktem Träger gilt.
  Mit Gleichungen (\ref{eq:limNablaP}) und (\ref{eq:convolutionNablaP}) erhält man
  $$
  \lim_{j \to \infty} \norm{\nabla p - \nabla \tilde p_j}_{L^2(\R^n)^n}
  \leq 
  \lim_{j \to \infty} \norm{\nabla p - \nabla p_j}_{L^2(\R^n)^n} + \lim_{j \to \infty} \norm{\nabla p_j - \nabla \tilde p_j}_{L^2(\R^n)^n}
  = 0, 
  $$
  woraus die Behauptung aus Gleichung (\ref{eq:GCharacterization}) folgt.

  Als Nächstes beweisen wir die Charakterisierungsgleichung (\ref{eq:LCharacterization}).
  Dazu sei im Folgenden
  $$
  L := \{ f \in L^2(\R^n)^n \mid \div f = 0 \}.
  $$

  Wir beweisen zunächst die Inklusion $L_\sigma^2(\R^n) \subseteq L$.
  Sei dazu $f \in L_\sigma^2(\R^n)$.
  Nach Definition existiert eine Folge $(f_j)_{j \in \N}$ mit $f_j \in C_{0,\sigma}^\infty(\R^n)$ für alle $j \in \N$ und 
  $$
  \lim_{j \to \infty} \norm{f_j - f}_{L^2(\R^n)^n} = 0.
  $$
  Daraus folgt unter Verwendung der \cauchy\hyp{}\schwartz\hyp{}Ungleichung
  \begin{align*}
    |[\div f, \varphi] |
    &= |[\div f_j - f, \varphi] | \\
    &= |\langle f_j - f, \nabla \varphi \rangle| \\
    & \leq \norm{f_j - f}_{L^2(\R^n)^n} \norm{\nabla \varphi}_{L^2(\R^n)^n} 
    \to 0, \quad j \to \infty,
  \end{align*}
  für alle $\varphi \in C_0^\infty(\R^n)^n$.
  Damit folgt $[\div f, \varphi] = 0$ für alle $\varphi \in C_0^\infty(\R^n)$.
  Es gilt also $\div f = 0$ im distributionellen Sinne und wir erhalten $f \in L$.

  Zum Beweis der umgekehrten Inklusion sei $f \in L$.
  Dann gilt
  $$
  \langle f, \nabla p \rangle = -[\div f, p] = 0
  $$
  für alle $p \in C_0^\infty(\R^n)$.
  Also ist die Abbildung $\langle f, \cdot\rangle$ stetig auf $\nabla (C_0^\infty(\R^n))$ und besitzt somit eine stetige Fortsetzung auf den Abschluss des Definitionsbereichs. 
  Dieser stimmt jedoch nach dem ersten Teil des Beweises mit $G(\R^n)$ überein.
  Daher erfüllt die Fortsetzung die Identität
  $$
  \langle f, \nabla p \rangle = 0 
  $$
  für alle $\nabla p \in G(\R^n)$.
  Dies bedeutet aber gerade, dass $f \in G(\R^n)^\perp$ gilt und aus Satz \ref{lem:helmholtz} zur \helmholtz\hyp{}Zerlegung folgt damit
  $$
  f \in G(\R^n)^\perp 
  = L_\sigma^2(\R^n)^{\perp\perp} 
  = \overline{L_\sigma^2(\R^n)}^{\norm{\cdot}_2}
  = L_\sigma^2(\R^n).
  $$
  Daraus folgt $L \subseteq L_\sigma^2(\R^n)$ und Gleichung (\ref{eq:LCharacterization}) ist bewiesen.
\end{proof}

\chapter{Zusammenfassung und Ausblick}

In dieser Arbeit wurde die Existenz und Eindeutigkeit \helmholtz\hyp{}Zerlegung des Funktionenraumes $L^2(\Omega)$ für ein Gebiet $\Omega \subseteq \R^n$ bewiesen.
Nach einer Wiederholung und Vertiefung integrationstheoretischer und funktionalanalytischer Grundlagen bestand die Hauptarbeit darin, unter geeigneten Zusatzannahmen die Lösbarkeit der im distributionellen Sinne zu verstehenden Gradientengleichung $\nabla p = f$ zu sichern.
Dazu war es zunächst nötig zu beweisen wie sich \lipschitz\hyp{}Gebiete geeignet von innen approximieren lassen.
Zudem beinhalteten die Hilfsresultate auch einen Darstellungssatz für Funktionale auf \sobolev\hyp{}Räumen.

Der Beweis der \helmholtz\hyp{}Zerlegung gestaltete sich aufgrund dieser Vorarbeit nunmehr übersichtlich.
Zusätzlich wurde eine Charakterisierung eines an der Zerlegung beteiligten Funktionenräume im Falle $\Omega = \R^n$ bewiesen.
In \cite[S.81ff., II.2.5]{sohr2001navier} finden sich weitere Charakterisierungen für beschränkte \lipschitz\hyp{}Gebiete sowie eine Dichtheitseigenschaft.

Anwendung findet die \helmholtz\hyp{}Zerlegung beispielsweise in der Konstruktion schwacher Lösungen der \stokes\hyp{}Gleichung \cite[S.129f.]{sohr2001navier}. 
Es lässt sich zum Beispiel zeigen, dass der \stokes\hyp{}Operator bis auf einen Vorfaktor mit der Verkettung der Operatoren $(P \circ \Delta)$ übereinstimmt, wobei $P$ die zur \helmholtz\hyp{}Zerlegung gehörige orthogonale Projektion und $\Delta$ den \laplace\hyp{}Operator darstellt.

Es stellt sich die Frage, inwieweit, sich die Existenz der \helmholtz\hyp{}Zerlegung und damit verbundener Aussagen auf weitere Funktionenräume verallgemeinern lassen.
Für allgemeine $L^q$\hyp{}Räume ist dies nicht möglich, wie Gegenbeispiele in \cite{maslennikova1986elliptic} oder \cite{bogovski1986decomposition} demonstrieren.
Es lässt sich jedoch zeigen, dass die Existenz der \helmholtz\hyp{}Zerlegung äquivalent zur Lösbarkeit eines geeigneten \neumann\hyp{}Problems ist. \cite{simader1992new}, \cite[S.146, Lemma III.1.2]{galdi2011navier}
Über diese Charakterisierung der Zerlegung wird in \cite[S.152, Theorem III.1.2]{galdi2011navier} für $C^2$-Gebiete oder Halbräume des $\R^n$ mit $n \geq 2$ die Existenz einer \helmholtz\hyp{}Zerlegung bewiesen.
Andererseits zeigen \cite{farwig05thehelmholtz,farwig05anLq} für beliebige Gebiete, dass sich die Resultate zur \helmholtz\hyp{}Zerlegung in \hilbert\hyp{}Räumen auf Schnitte beziehungsweise direkte Summen bestehend aus einem $L^2$- und einem $L^q$\hyp{}Raum übertragen lassen.



\cleardoublepage
%\nocite{*}
\bibliographystyle{alphadin}
\bibliography{bachelorthesis_fabian_gabel}
%\addcontentsline{toc}{chapter}{Literaturverzeichnis}
\addcontentsline{toc}{chapter}{Literaturverzeichnis}


%\cleardoublepage
%\renewcommand{\glossarysection}[2][\theglstoctitle]{%
%    \def\theglstoctitle{#2}%
%      \setcounter{secnumdepth}{-1}%
%        \chapter[#1]{#2}%
%}
%\printglossary[title=Symbolverzeichnis,toctitle=Symbolverzeichnis]
      
\cleardoublepage
% Keine Kopf-/Fußzeile auf dieser Seite
\thispagestyle{empty}

\vspace*{4cm}
\section*{Erklärung}\index{Erklärung}

Hiermit versichere ich, dass ich die vorliegende Arbeit selbständig verfasst habe und alle benutzten Quellen einschließlich der Quellen aus dem Internet und alle sonstigen Hilfsmittel angegeben habe.\vspace{20pt}

\noindent
Darmstadt, den \today\vspace{50pt}

% Unterschrift (handgeschrieben)

\noindent
Fabian Gabel


\end{document}
